\documentclass[12pt, a4paper]{article}
\usepackage{graphicx}
%\usepackage[a4paper,margin=1in]{geometry}
\usepackage{hyperref}
\usepackage{amsmath}
\usepackage{amssymb}
\usepackage{natbib} 
\usepackage{siunitx}

\title{Molecular Docking Programs}
\author{Elvis Martis}
\begin{document}
\maketitle
%\chapter*{Molecular Docking Methods and Programs}
\section{Introduction}
As structure-based drug design developed into a mature science, molecular docking got better with every step. The method depends on how well the conformational search algorithm can find the bioactive conformation(s). However, the search algorithm by itself does not know what a bioactive conformation(s) looks like. By definition, bioactive conformation is the atomic arrangement of a molecule that is suitable for binding to its respective target. It either inhibits/antagonises or activates it. Practically speaking, the bioactive conformation is always one of the plausible low-energy states, not necessarily the global minimum. Furthermore, there is always a possibility that a ligand may modulate more than one target, and thus, the bioactive conformations may differ for the same ligand based on its target receptor. Remember that the conformational search tool searches for the bioactive conformation within the receptor binding site; therefore, many low-energy conformations that are observed in the gas phase or in solution phase may not be sampled due to clashes with the receptor atoms. Therefore, docking programs implement a preliminary "bump-check" that can filter out conformations that are likely to show clashes with the receptor atoms. This way, more intense energy calculations are performed on fewer conformationsm, which improves the speed of docking. 

Now that we have established that there are several possible bioactive conformations for a ligand depending upon its binding target, the conformational search algorithm alone is unable to identify it, but it can find it among the pool of low energy states. The scoring function assigns a numerical value to the each binding pose, which depicts the binding affinity. This almost always is not the free energy of binding ($\Delta$G$_{\text{bind}}$), but merely an estimate of binding enthalpy ($\Delta$H$_{\text{bind}}$). Based on this affinity value, called as \textit{docking score}, the scoring functions ranks all the poses based on what it finds most like conformation to bind to the target receptor. The chapter on scoring functions (\textbf{Link to chapter}) deals with this part in detail. 


With significant improvements in computing power, and ease of computing led to the development of several molecular docking programs. Understandably, it is impossible to track and understand every bit of the development. Therefore, to assist, particularly, the beginners, this chapter will describe few popular molecular docking programs (\textbf{Table \ref{tab:List}}).  Knowing every program to make a choice that best the problem at hand is almost always impossible. However, this is also not necessary as there exists ample literature reporting the use of many of the tools with their outcome well-documented. Irrespective of the choice of the docking program, the general steps will remain same; for example, preparation of the receptor, defining the binding site as search space, preparation of ligands, specifying the docking parameters, validating the docking procedure, performing the docking calculations, and analysis and dissemination of the results.  In a nutshell, the general principles remain unchanged; however, there are differences in the file format of the ligand, docking and scoring algorithms. Most of this information is well-documented in the user manual. it is strongly recommend that, before using any docking algorithm (in fact true for all computational tools/software), the user manual must be consulted to understand the importance and consequence of each parameter. The developers present their report on several validation studies as part of the software development, which is a great resource to understand the effect of each parameter on the outcome. A literature survey of the program's usage will also be covered in the user manual. 

\begin{table}[ht]
    \centering
    \caption{Non-exhaustive list of common molecular docking software}
    \label{tab:List}
    \begin{tabular}{|c c c |}
    \hline
         \textbf{Software} & \textbf{License type} & \textbf{Scoring Function}  \\
    \hline
        AutoDock4 & Open Source & Force field  \\
        AutoDock-Vina & Open Source & Empirical  \\
        DOCK & Open Source (Academic) & Force field  \\
        GOLD &  Commercial & Multiple options \\
        FlexX &  Commercial & Empirical \\
        Glide &  Commercial & Empirical \\
        FlexX & Commercial & Multiple options \\
        SurFlex & Commercial & Multiple options \\
    \hline
    \end{tabular}
\end{table}


\section{DOCK}
The DOCK program is considered the first tool developed for molecular docking in 1982\cite{CH7_Kuntz1982}. It was developed by the University of California, San Francisco (UCSF) in the group of Irwin Kuntz. It was the first automated procedure to dock small molecules into a receptor site of known structure, and remains under active development in the form of DOCK3 and DOCK6 variants\cite{CH7_Meng1992Grid,CH7_Allen2015,CH7_Brozell2012}. It aims to predict the most favourable binding conformation of small molecules within the active sites of macromolecular targets, such as proteins. The first version of DOCK could only tackle rigid-body docking, which means the receptor and ligand both are treated as rigid entities with only possibility to rotate and orient the ligand as a whole without rotating the torsion angles to explore the ligand conformations in the binding site. In such cases, an external conformation search tool is used to explore the ligand conformations which are docked as rigid bodies. DOCK then uses a geometric matching algorithm that aligns the ligand with a negative representation of the receptor's binding pocket. With several improvements single the first version, and increase in the computational power, its functionality has evolved to incorporate features such as force field based scoring, real time optimisation, and flexible ligand docking. 

The DOCK programs accommodates various input file formats, including PDB, MOL2, and SDF, and is extensively utilised for virtual screening and lead optimisation in drug development. Nearly, all docking programs will break down the docking step in to two fundamental process as shown below: 
\begin{enumerate}
  \item \textbf{Sampling}: Generating ligand conformations and orientations in the binding site.
  \item \textbf{Scoring}: Evaluating each conformation (pose) using scoring functions that evaluates inter-molecular interactions.
\end{enumerate}


\subsection{The theory underlying DOCK program}

\subsubsection{Binding-site representation: spheres and shape complementarity}


The DOCK program represents potential binding sites by filling surface invaginations with a set of overlapping spheres, effectively creating a negative image of the receptor’s topography using an accessory program called sphgen\cite{CH7_Kuntz1982}. 
\begin{itemize}

\item \textit{Generation Process}: The process begins by calculating the receptor's molecular surface defined by rolling a probe the size of a water molecule over the van der Waals surface. The sphgen program then generates spheres analytically: each sphere touches the molecular surface at two points and has its center located along the surface normal of one of those points.
\item \textit{Spatial Orientation}: For the receptor, spheres are placed on the outside of the surface to fill pockets and grooves. For the ligand, spheres are placed on the inside of its surface to approximately fill the space it occupies.
\item \textit{Filtering and Accuracy}: To manage computational complexity, DOCK filters these spheres, typically retaining only the largest sphere per surface atom that does not intersect the molecular surface. The radius of these spheres reflects local concavity; larger radii describe shallower pockets, while smaller radii describe deeper invaginations.
\item \textit{Clustering}: Overlapping spheres are grouped into clusters using a single-linkage algorithm. Each cluster represents a presumptive binding site, and in most examined proteins, the largest cluster corresponds to the actual binding site. 
\end{itemize}

\subsubsection{Ligand–sphere matching and pose generation}

The matching process is based on internal distance comparisons between the centers of the receptor spheres and ligand spheres\cite{CH7_Kuntz1982,CH7_Shoichet1992Shape}:

\begin{enumerate}

  \item \textit{The Pairing Rule}: A small subset (typically 3 or more) of sphere centers in the binding site are choosen. A ligand sphere set can be paired with a receptor sphere set if all pairwise internal distances between the ligand spheres match the corresponding distances between receptor spheres within a specified error tolerance \(\varepsilon\):
        \[
          \left| d_{PiPj} - d_{LiLj} \right| \le \varepsilon,
        \]
        where $d_{PiPj}$ is the distance between sphere centers Pi and Pj, and$(d_{LiLj}$ is the distance between ligand atoms Li and Lj.
        
        \item \textit{Graph Theoretical Approach}: The problem is treated as a bipartite graph. Nodes represent potential pairings between ligand and receptor descriptors; an edge exists between two nodes if the distance between the two ligand descriptors matches the distance between the two receptor descriptors.
        
        \item \textit{Tree Pruning}: To avoid a combinatorial explosion, DOCK uses a build-up procedure. It evaluates the graph as each node is added; if a new node fails the distance check with previously assigned pairs, that entire branch of the search tree is "pruned".
      
       \item \textit{Search Algorithms}: Later versions of DOCK introduced three specific methods for constructing these matching graphs:
       
       \begin{itemize}
       \item Fan Algorithm: Selects descriptors based on their distance from an initial starting point in each ligand.
        \item Cat's Cradle Algorithm: Selects descriptors at each level based on their distance to the descriptor successfully used at the previous level.
        \item Center of Mass Algorithm: Selects descriptors based on their distance to the ligand's center of mass.
       \end{itemize}

\end{enumerate}

Once a valid match of spheres is identified, DOCK generates a specific orientation (pose) for the ligand.
\begin{enumerate}
\item \textit{Least-Squares Superimposition}: A minimum of four nonplanar sphere pairs is required to uniquely define a configuration. DOCK uses a least-squares algorithm (such as ORIENT\cite{CH7_ORIENT}) to calculate a rotation/translation matrix that superimposes the ligand spheres onto their matched receptor sphere partners.
\item \textit{Coordinate Transformation}: This transformation matrix is applied to all atoms of the ligand to generate its  structure within the receptor's coordinate system.
\item \textit{Refinement and Filtering}: The resulting poses are subjected to a fast evaluation. Any orientation that results in atomic overlap (placing the ligand in space already occupied by the receptor) is discarded. The program may also perform a local optimization by moving ligand atoms toward the centers of the closest receptor spheres to improve the fit.
\item \textit{Output}: The final output is a set of geometrically feasible ligand structures, each scored based on the degree of overlap or other complementarity metrics
\end{enumerate}

In early versions, DOCK treated ligands as rigid body. Subsequently, flexible docking strategies such as anchor-and-grow and genetic algorithms were introduced, which treat the ligand as a set of rotatable fragments and sample torsional degrees of freedom during or after initial rigid-body placement\cite{CH7_Oshiro1995,CH7_Meng1993Min}.

\subsubsection{Grid-based interaction energy evaluation}

Once a pose is generated, DOCK typically evaluates its quality using precomputed interaction grids. A separate program grid, computes the receptor contribution to a molecular mechanics-type potential on a 3D lattice surrounding the binding site \cite{CH7_Meng1992Grid}. The interaction energy between ligand and receptor is approximated as the sum of van der Waals and electrostatic terms:
\[
  E_{\text{int}} \approx E_{\text{vdw}} + E_{\text{es}}.
\]

The key idea is to precompute the receptor-dependent part of the interaction energy at each grid point k and store it in grid files. During docking, these values are interpolated and combined with ligand-dependent parameters to yield a fast approximate energy evaluation for each pose. This allows significant speeds during docking and scoring because receptor
contributions are computed only once which are generally compute intensive due to size of the receptor compared to the ligand. While the ligand placements are evaluated many thousands to millions of times.

\subsubsection{Approximate nature of the scoring function}

The inter-molecular interaction energy is modeled as a sum of:
\begin{itemize}
  \item Lennard–Jones-type van der Waals interactions
  \item Coulombic electrostatics, often with a distance-dependent dielectric
\end{itemize}

It typically does not include explicit hydrogen bonding, solvation, or conformational entropy terms in the primary scoring function\cite{CH7_Meng1992Grid}. However, with later versions of DOCK program more sophisticated scoring and rescoring functions (e.g.Poisson Boltzmann/Surface Area (Delphi), Generalised Born/Surface Area, AMBER-based scoring (Chemgrid),
similarity) are available \cite{CH7_Allen2015,CH7_Brozell2012}. 


\subsection{A typical DOCK Workflow}

The DOCK program offers a highly configurable workflow. A typical DOCK6 workflow that is suitable structure-based virtual screening involves the following steps\cite{CH7_Allen2015}.

\subsubsection{Receptor preparation}

\begin{enumerate}
  \item Start from an experimental (X-ray, cryo-EM, or NMR) or computed structure models (AlphaFold, Rosetta etc)  of the protein. It recommended to use a structure with a bound ligand which makes the binding site suitable for docking.
    \item Cleaning the structure involves following steps:
        \begin{itemize}
          \item Remove crystallographic waters unless known to be functionally important.
          \item Add missing side chains or loops if required.
          \item Decide on protonation/tautomeric states (pH-dependent).
        \end{itemize}

  \item Convert to a format with atom types and partial charges
        (typically MOL2), using external tools (e.g. UCSF Chimera, which is recommended).
  \item Define the binding site (e.g. around a bound ligand or specified residues) and create a bounding box (via
        showbox or similar utilities).
\end{enumerate}

\subsubsection{Sphere generation (sphgen) and clustering}

\begin{enumerate}
  \item Use sphgen to fill the binding site with overlapping spheres that approximate the shape of the pocket:
  \item Optionally, cluster and prune spheres using sphcluster to reduce the redundancy or focus on a particular sub-pocket
  \item Optionally, color spheres with simple chemical labels (donor/acceptor/hydrophobe, etc.) for chemical matching.
\end{enumerate}

\subsubsection{Grid generation (grid)}

\begin{enumerate}
  \item Define a 3D grid box that encloses the binding site (often oriented to the sphere cluster defined above).
  \item Run grid program to precompute the following grid maps:
        \begin{itemize}
          \item Bump grid (*.bmp file), which encodes severe steric clashes.
          \item Contact grid (*.cnt file), which is used for contact-based (shape-like) scoring.
          \item Energy grid (*.nrg file), which is used for force-field-like scoring (van der Waals + electrostatics).
        \end{itemize}

  \item Choose parameters such as:
        \begin{itemize}
          \item Grid spacing (commonly 0.3–0.5 \r{A}).
          \item Lennard–Jones exponents (e.g.\ 12–6 or 9–6).
          \item Dielectric model (distance-dependent vs. constant).
          \item Cutoff distances for van der Waals and electrostatic contributions.
        \end{itemize}
\end{enumerate}

\subsubsection{Ligand preparation}
The ligand preparation is more or less done in a similar manner for almost all docking programs, therefore, a chapter is fully dedicated to small molecular handling and preparation (\textbf{link to chapter}). Most programs differ in the methods employed for partial charge calculations, atom-typing and output file format.

\begin{enumerate}
  \item Generate 3D coordinates for the ligands (e.g. from SMILES or a database).
  \item Enumerate relevant protonation and tautomeric states.
  \item Assign atom types and partial charges compatible with the receptor parameterization.
  \item Store ligands in MOL2 format suitable for DOCK program.
\end{enumerate}

\subsubsection{Docking and scoring}

\begin{enumerate}
  \item Run the main docking program (executed using \textit{dock6} command) with a chosen sampling strategy:
        \begin{itemize}
          \item Rigid-body docking using sphere–atom matching.
          \item Flexible docking using anchor-and-grow, genetic algorithms, or other methods.
        \end{itemize}

  \item For each generated pose:
        \begin{enumerate}
          \item Apply bump checking to discard severe clashes.
          \item Evaluate one or more scoring functions:
                \begin{itemize}
                  \item Grid-based energy (van der Waals + electrostatics).
                  \item Contact (shape-like) score.
                  \item Continuous Cartesian energy (non-grid).
                  \item Footprint similarity or pharmacophore scores.
                \end{itemize}

          \item One may also perform a local minimization (rigid-body or torsional) to relax the pose into a nearby local minimum.
        \end{enumerate}
  \item Store and rank poses based on primary scores. Also, note there was a secondary score which would score initial ligand conformations, however, this is now obsolete from DOCK version 6.10 onwards.
\end{enumerate}

\subsubsection{Post-processing}

\begin{enumerate}
  \item Cluster poses by RMSD to identify distinct binding modes.
  \item Visually inspect top-ranked poses for key interactions, and steric clashes
  \item Further, rescore with more sophisticated methods (PB/SA, GB/SA, AMBER score, footprint similarity, etc.).
\end{enumerate}


\subsection{Scoring Functions in DOCK}

DOCK provides several scoring functions, many of which can be combined through the descriptor score (also called consensus scoring) framework in DOCK6\cite{CH7_Allen2015}. 
The following are the scoring functions available in DOCK:
\begin{itemize}
  \item Force field-based grid energy score (\textit{grid score} or \textit{energy score}).
  \item Contact (shape-like) score.
  \item Continuous (Cartesian) energy score.
  \item Footprint similarity and pharmacophore-based scores (for more advanced workflows).
\end{itemize}

\subsubsection{Force field-based grid energy score}

\begin{enumerate}
\item \textit{Pairwise interaction}
The starting point is a pairwise additive interaction potential between ligand atoms $i \in L$ and receptor atoms
$j \in R$:

\[
  E_{\text{int}} = E_{\text{vdw}} + E_{\text{es}},
\]
with
\[
  E_{\text{vdw}} =
    \sum_{i \in L} \sum_{j \in R}
    \left(
      \frac{A_{ij}}{r_{ij}^{n}} -
      \frac{B_{ij}}{r_{ij}^{m}}
    \right),
\]
\[
  E_{\text{es}} =
    \sum_{i \in L} \sum_{j \in R}
    \frac{q_i q_j}{\varepsilon(r_{ij})\, r_{ij}},
\]
where:
\begin{itemize}
  \item $r_{ij}$ is the distance between ligand atom i and receptor atom j.
  \item $A_{ij}$ and $B_{ij}$ are Lennard--Jones parameters derived from atomic van der Waals (vdW) radii and well depths
        (via combining rules).
  \item n and m are the exponents of the generalized Lennard--Jones potential (commonly n = 12, m = 6; vdW soft potentials can be used by defining n = 9, m = 6).
  \item $q_i$ and $q_j$ are partial charges of ligand and receptor atoms, respectively.
  \item \(\varepsilon(r_{ij})\) is the effective dielectric chosen as a distance-dependent form, e.g.
        \(\varepsilon(r) = \varepsilon_0\, r\) with
        \(\varepsilon_0 \approx 4\).
\end{itemize}

The DOCK programs allows for user-defined scaling these  energetic contributions as shown here:
\[
  E_{\text{grid}} =
    w_{\text{vdw}}\, E_{\text{vdw}} +
    w_{\text{es}}\, E_{\text{es}},
\]
where $w_{\text{vdw}}$ and $w_{\text{es}}$ are user-defined scaling factors.

\item \textit{Generalized Lennard--Jones form}

For identical atoms with vdW radius R and well depth\(\varepsilon\), DOCK employs a generalized Lennard--Jones expression:
 \[
  V_{\text{LJ}}(r) = \varepsilon
    \left[
      C \left( \frac{R}{r} \right)^{n}
      -
      D \left( \frac{R}{r} \right)^{m}
    \right],
\]
with exponents $n > m$ (e.g. n = 12, m = 6 or n = 9, m = 6).  Constants C and D are chosen such that:

\begin{align*}
  V_{\text{LJ}}(r = R) &= -\varepsilon, \\
  \frac{\mathrm{d} V_{\text{LJ}}}{\mathrm{d}r}\biggr|_{r = R} &= 0,
\end{align*}
which fixes the position and depth of the potential minimum while allowing the repulsive wall to be softened by decreasing the 12--6 potentials to 9--6 potentials).

\item \textit{Grid-based formulation}
Computing the full double sum over all receptor atoms is computational intense, therefore, DOCK precomputes the receptor contribution on the grid maps \cite{CH7_Meng1992Grid}. For each grid point k, it stores three receptor-dependent
quantities:
\[
  G^{(n)}_k = \sum_{j \in R}
    \frac{A_j}{r_{jk}^{n}},
  \qquad
  G^{(m)}_k = \sum_{j \in R}
    \frac{B_j}{r_{jk}^{m}},
  \qquad
  G^{(\text{es})}_k =
    \sum_{j \in R}
    \frac{q_j}{\varepsilon(r_{jk})\, r_{jk}},
\]
where $r_{jk}$ is the distance between receptor atom j and grid point k, and $A_j$ and $B_j$ are atom only Lennard--Jones
parameters (before combining with ligand atoms). 

Given a ligand pose, each ligand atom i is mapped to the surrounding grid points (via nearest-neighbor or trilinear
interpolation). Denoting the effective grid index for atom i by $k(i)$, the interaction energy is computed as:
\[
  E_{\text{vdw}} \approx
    \sum_{i \in L}
      \left(
        A_i\, G^{(n)}_{k(i)} -
        B_i\, G^{(m)}_{k(i)}
      \right),
\]
\[
  E_{\text{es}} \approx
    \sum_{i \in L}
      q_i\, G^{(\text{es})}_{k(i)}.
\]
Here $A_i$ and $B_i$ are ligand atomic parameters, and the precomputed receptor sums $G^{(n)}_k$, $G^{(m)}_k$, and
$G^{(\text{es})}_k$ are read from the energy grid files. This allows to reduce the computational complexity from $O(N_L N_R)$ double sum into $O(N_L)$ sum over ligand atoms, where $N_L$ and $N_R$ are the number of ligand and receptor atoms, respectively.

\item \textit{Contact (shape-like) scoring}

The contact score is a simple, shape oriented scoring function that counts favorable close contacts between ligand and receptor heavy atoms, while penalizing severe overlaps. The contact score emphasizes steric complementarity and is often used in combination with other more physically detailed scores. A contact is defined by a distance cutoff
$d_{\text{cut}}$, typically around 4.5 \r{A}, and a bump is defined by an overlap factor $\alpha$ of the sum of vdW radii:

\begin{itemize}
  \item If
        $\alpha (R_i + R_j) \le r_{ij} \le d_{\text{cut}}$, atoms i and j are considered to be in contact.
  \item If
        $r_{ij} < \alpha (R_i + R_j)$, the atoms are in severe steric clash (a bump).
\end{itemize}

A simple conceptual form of the contact score is:
\[
  S_{\text{contact}} =
    \sum_{i \in L} \sum_{j \in R} f(r_{ij}),
\]
with a pairwise function
\[
  f(r_{ij}) =
  \begin{cases}
    s_{\text{bump}} &
      \text{if } r_{ij} < \alpha (R_i + R_j), \\
    s_{\text{contact}} &
      \text{if } \alpha (R_i + R_j) \le r_{ij} \le d_{\text{cut}}, \\
    0 & \text{otherwise},
  \end{cases}
\]
where:
\begin{itemize}
  \item \(s_{\text{contact}} < 0\) denotes a favorable contact contribution.
  \item \(s_{\text{bump}} > 0\) is a penalty for steric clash.
\end{itemize}


\item \textit{Continuous (Cartesian) energy score}

Instead of relying on grids, DOCK can also compute the same non-bonded interaction energy in continuous Cartesian space. Continuous scoring avoids grid discretization artifacts and can be used as a more accurate secondary score or during local minimization, at the cost of higher computational expense.
This continuous score uses the same functional form as the grid-based energy:
\[
  E_{\text{cont}} =
    w_{\text{vdw}}\, E_{\text{vdw}} +
    w_{\text{es}}\, E_{\text{es}},
\]
but:
\begin{itemize}
  \item Distances $r_{ij}$ are computed directly between ligand and receptor atoms (no grids).
  \item Lennard--Jones exponents and dielectric models are adjustable (e.g. 12--6 or 9--6; distance-dependent or constant dielectric).
\end{itemize}


\item \textit{Footprint similarity and pharmacophore-based scores}

Modern DOCK6 versions provide more elaborate scoring functions, often used in combination via the descriptor score framework\cite{CH7_Allen2015}. These advanced scores are particularly useful when one wants to enforce similarity to a known active ligand or refine virtual screening hits. 

\begin{itemize}
  \item Footprint similarity score: decomposes interactions (van der Waals, electrostatics, hydrogen bonding) per receptor residue and compares the resulting footprint vectors between a reference and a candidate ligand using metrics such as Euclidean distance or Pearson correlation.
  \item MultiGrid score: computes residue-decomposed interaction energies on multiple grids and combines them with footprint similarity.
  \item Pharmacophore matching similarity: evaluates the geometric overlap between pharmacophore feature sets for a reference and a candidate.
  \item Descriptor score: a linear combination of multiple component scores:
        \[
          S_{\text{descriptor}} =
            \sum_k w_k\, S_k,
        \]
        where $S_k$ may be grid score, continuous score, footprint similarity, pharmacophore similarity, or other implemented
        descriptors.
\end{itemize}
\end{enumerate}


\section{GLIDE}

Glide (Grid-based Ligand Docking with Energetics) is a widely used molecular docking program developed by Schr\"{o}dinger\cite{CH7_Friesner2004a,CH7_Friesner2004b}. It combines a hierarchical search of ligand poses in a rigid receptor binding site with empirical, grid-based scoring functions derived from force field and knowledge-based terms. In its standard-precision (SP) and extra-precision (XP) modes, Glide has been shown to give accurate pose prediction and strong enrichment in retrospective screening benchmarks \cite{CH7_Friesner2004a,CH7_Friesner2006,CH7_Friesner2004b}.


\subsection{The theory underlying GLIDE program}
Glide works with the following assumption to improve the speed and accuracy of the docking results:
\begin{itemize}
  \item A rigid receptor (aside from optional scaling of van der Waals radii or induced-fit workflows handled outside the core Glide engine).
  \item Internally flexible ligands, with conformations generated and optimized on-the-fly or supplied externally (using confgen\cite{CH7_confgen}).
  \item A precomputed receptor grid that encodes nonbonded interaction fields (van der Waals, electrostatic, hydrophobic, etc.) on a 3D lattice.
\end{itemize}

The search algorithm is hierarchical and funnel shaped \cite{CH7_Friesner2004a}. Only the best-scoring refined poses, typically on the order of 100–400 per ligand, are passed to the final stages.

The key stages are given as follows:
\begin{enumerate}
  \item \textit{Stage 1 Site-point search}: The binding site is discretized into site points on a coarse 2 \r{A} grid covering the active-site region. For each ligand core conformation, distances from site points to the receptor surface are compared to distances from the ligand center to its molecular surface. Only those site points whose local shape is compatible with the ligand are retained as candidate placements.

  \item \textit{Stage 2 Greedy scoring and refinement}: For each retained placement, a series of filters is applied:
    \begin{itemize}
      \item Stage 2a Diameter test: Selected orientations of the ligand diameter (line through the two most distant atoms) are tested, discarding poses with severe steric clashes.
      \item Stage 2b Subset test: A subset of polar/metal-binding atoms is scored against the receptor using a discretized, ChemScore-like function.
      \item Stage 2c Greedy scoring: All ligand atoms are scored with a discretized empirical potential patterned after ChemScore, rewarding hydrophobic contacts, hydrogen bonds, and metal interactions, while penalizing steric clashes.
      \item Stage 2d Refinement: Top poses from greedy scoring are rigidly refined by small Cartesian translations to improve complementarity.
    \end{itemize}

  \item \textit{ Stage 3 Grid-based energy minimization}: The poses that survive are minimized on precomputed OPLS-AA van der Waals and Coulomb grids for the receptor. To improve numerical stability, minimization typically starts on smoothed grids and then anneals onto the full OPLS-AA nonbonded surface.

  For internally generated conformations, this stage includes:
    \begin{itemize}
      \item Rigid-body translation and rotation of the ligand.
      \item Torsional optimization of core and terminal rotamer groups.
    \end{itemize}

  \item \textit{Stage 4 Final scoring}: The minimized poses are evaluated with Glide's proprietary empirical scoring function, GlideScore SP (for standard precision) or GlideScore XP (for extra-precision mode). In case, the user wishes to perform virtual screening on a large chemical database, it is recommended to filter unwanted ligands using GlideScore HTVS (high-throughput virtual screening mode). Pose selection within a ligand and ranking across ligands use a combination of GlideScore and a model energy, Emodel, that places additional weight on the underlying nonbonded interaction energy, and the internal strain energy of the ligand conformation (for flexible docking). It is must noted that Emodel is appropriate to select the best pose amongst various docking poses of the same ligand, which Glidescore is more suitable to select best binders amongst various ligands. 
\end{enumerate}

By optimizing rotatable groups one at a time for a given core placement, Glide avoids explicit enumeration of all torsional combinations. For example, if a ligand has five rotamer groups with three discrete states each, a naïve exhaustive search would require evaluation of $3^5 = 243$ conformers, whereas sequential optimization of each group reduces this to $3 \times 5 = 15$ configurations. This decomposition, together with the hierarchical filters, makes large-scale virtual screening tractable while retaining good pose accuracy.

\subsection{Typical GLIDE Workflow}

The Glide module is bundled with Schr\"{o}dinger's Maestro graphical user interface. It provides a user friendly interface with all modules provided by  Schr\"{o}dinger Suite are integrated, thus, workflows can be easily implemented for all kinds drug design and discovery projects. 

\subsubsection{Protein and Ligand Preparation}

\begin{enumerate}
  \item \textit{Protein preparation}: The input receptor structure is cleaned using Schr\"{o}dinger's Protein Preparation Wizard (PPW). As corrective measures, the PPW performs following operations on the input receptor structure:
    \begin{itemize}
      \item Adding hydrogens and assigning bond orders.
      \item Determining protonation states of ionizable groups at the target pH using Epik\cite{CH7_epik} module.
      \item Optimising the protein interaction network by flipping Asn/Gln/His as needed.
      \item Missing side chains are added using the prime module.
      \item Finally, minimizing the structure with restrained heavy-atom coordinates, thus sresolving steric clashes.
    \end{itemize}

  \item \textit{Ligand preparation}: Ligands are prepared with LigPrep and/or Epik:
    \begin{itemize}
      \item 3D coordinate generation with low-energy conformations using the LigPrep module.
      \item Generation of tautomers and protonation states over a desired pH range using the Epik module.
      \item Standardization of atom types, formal charges, stereochemistry, and ring conformations.
    \end{itemize}
\end{enumerate}

\subsubsection{Receptor Grid Generation}

Glide requires a receptor grid that encodes receptor shape and interaction fields on a 3D lattice. The idea grid-based docking is similar to previously explained DOCK program, which is essential for docking speed ups.
The key steps are as follows:
\begin{itemize}
  \item \textit{Defining the binding site}: The grid center is defined based on a bound ligand, a set of active-site residues, or a region identified by a binding site detection tool (SiteMap\cite{CH7_sitemap}).
  \item \textit{Box dimensions}: An inner box defines the allowed ligand centroid region, and an outer grid box defines the full volume sampled by ligand atoms.
  \item \textit{Nonbonded interaction fields}: OPLS-AA van der Waals and Coulomb potentials are computed on the grid, often with scaled radii/charges to allow limited receptor softness.
  \item \textit{Optioanal constraints}: Hydrogen-bond, metal, and hydrophobic constraints can be encoded in the grid to enforce key pharmacophoric interactions during docking.
\end{itemize}

\subsubsection{Docking Settings and Modes}

Glide offers several precision levels with increasingly severe penalty terms, which is known to enhance discrimination between binders and non-binders at an additional computational cost. In a typical workflow, large libraries are docked first with HTVS or SP, and only the top fraction (e.g., 5--10\%) is redocked with XP for more rigorous scoring.

\begin{itemize}
  \item HTVS (High-Throughput Virtual Screening): Aggressive use of filters and simplified scoring for very fast screening of large libraries. Sampling is reduced; intended for initial triage rather than final pose analysis.
  \item SP (Standard Precision): Default mode for most docking tasks. Uses the full hierarchical search and SP GlideScore, providing a balance between speed and accuracy. Suitable for pose prediction and primary ranking of hits \cite{CH7_Friesner2004a,CH7_Friesner2004b}.
  \item XP (Extra Precision): It is a rigorous, computationally intensive scoring function designed to accurately estimate protein-ligand binding affinities while primarily focusing on minimising false positives. Unlike the "softer" Standard Precision (SP) mode, XP GlideScore is a "hard" function that exacts severe penalties for violations of physical chemistry principles, such as burying polar groups without adequate solvation \cite{CH7_Friesner2006}.
\end{itemize}


\subsubsection{Pose Selection and Post-processing}

It is always recommended that any docking output must be evaluated for multiple docking poses for the same ligand and not merely the top/best pose. It is never guarenteed that top pose is always the physiological relevant pose, despite the fact the all developers aim to achieve this. The post-processing step includes, but not limited to, visual inspection, clustering of poses, rescoring (e.g., with prime/MM-GBSA), and comparison with experimental SAR. 

\begin{itemize}
  \item \textit{Emodel}: An internal model energy that combines GlideScore, the nonbonded grid energy (Coulomb + van der Waals), and, for flexible docking, the ligand's internal strain energy. Emodel emphasizes the physical interaction energy and is primarily used to select the best pose of each ligand.
  \item \textit{GlideScore}: An empirical scoring function estimated the approximate free energy binding. It is optimized for docking accuracy and enrichment. GlideScore is used to compare different ligands; more negative values indicate more favorable predicted binding \cite{CH7_Friesner2004b}.
\end{itemize}


\subsection{Standard-Precision GlideScore: Functional Form and Terms}

In SP (and related HTVS) docking, the default scoring function is GlideScore, an empirical, ChemScore-inspired function \cite{CH7_Friesner2004a,CH7_Friesner2004b}. GlideScore thus blends force field-based interaction energies ($\text{E}_{\text{vdW}}$, $\text{E}_{\text{Coul}}$) with empirically fitted terms capturing hydrophobic enclosure, hydrogen bonding, metal coordination, desolvation penalties, and ligand conformational entropy. The functional form and coefficients are optimized against experimental binding affinities and database enrichment metrics across diverse targets \cite{CH7_Friesner2004b,CH7_Friesner2006}. We would like to point out that the coefficients mentioned in the following equations may change with newer version of Glide scoring functions due to constant developments and improvements made. 


The functional form of SP GlideScore (\(G_{\text{Score}}\)) is given by:
\[
  G_{\text{Score}} \;=\; 0.065\,E_{\text{vdW}} \;+\; 0.130\,E_{\text{Coul}} \;+\; E_{\text{Lipo}} \;+\; E_{\text{HBond}} \;+\; E_{\text{Metal}} \;+\; E_{\text{BuryP}} \;+\; E_{\text{RotB}} \;+\; E_{\text{Site}}
\]
where all terms are in units of kcal/mol. By convention, more negative values of $\text{G}_{\text{Score}}$ correspond to stronger predicted binding affinity. 

The components and coefficients are as follows
\begin{enumerate}
  \item $\text{E}_{\text{vdW}}$ (vdW): Van der Waals interaction energy between ligand and receptor atoms is evaluated on the receptor grid with a reduced net ionic charges on groups with formal charges (e.g., metals, carboxylates, guanidinium groups). The coefficient 0.065 scales this term to an empirical binding energy like contribution. 

  \item $\text{E}_{\text{Coul}}$ (Coul): Coulombic (electrostatic) interaction energy between ligand and receptor atoms are also computed on the receptor grid with reduced formal charges. The coefficient 0.130 reflects a somewhat larger relative contribution of electrostatics in the calibrated scoring function.

  \item $\text{E}_{\text{Lipo}}$ (Lipo): Lipophilic/hydrophobic term derived from a hydrophobic grid potential that rewards favorable burial of nonpolar surface area and lipophilic–lipophilic contacts. This term captures the hydrophobic effect in an empirical fashion, complementing the van der Waals term.

  \item $\text{E}_{\text{HBond}}$ (HBond): Hydrogen-bonding term is decomposed internally into components for:
    \begin{itemize}
      \item Neutral–neutral donor–acceptor pairs.
      \item Neutral–charged pairs.
      \item Charged–charged pairs.
    \end{itemize}

    Each component carries different weights and distance/angle dependence, allowing stronger reward for well-oriented, partially buried hydrogen bonds and lesser reward for solvent-exposed or suboptimal interactions.

  \item $\text{E}_{\text{Metal}}$ (Metal): Metal-binding term considers only a interaction between metal ions and anionic ligand acceptor atoms if present. If the net metal charge in the apo protein is positive, the scoring function includes a preference for anionic ligands; if the net charge is zero, this preference is suppressed.

  \item $\text{E}_{\text{BuryP}}$ (BuryP): The burial penalty term strongly penalises burying the polar groups, either on ligand or receptor, that are not engaged in satisfactory hydrogen bonding or ionic interactions. This term ensures that such physically unrealistic poses where charged or strongly polar atoms are desolvated without compensating interactions \cite{CH7_Friesner2004b,CH7_Friesner2006} do not rank high. 

  \item $\text{E}_{\text{RotB}}$ (RotB): Even though entropy term in not explicitly included in the scoring functions, this penalty term implicitly attempts added entropy penalty for freezing rotatable bonds upon binding.  More rotatable bonds typically incur a larger penalty unless offset by strong favorable interactions. 

  \item $\text{E}_{\text{Site}}$ (Site): If polar groups, not involved in hydrogen bonding, are placed in polar microenvironment of an otherwise non-polar site, it rewards such a placement, otherwise a penalty is added to discourage incorrect positioning of such groups. 

\end{enumerate}


\subsection{XP GlideScore: Extended Model of Binding}

The XP (extra-precision) scoring function introduces a more detailed description of protein–ligand recognition and desolvation, especially for hydrophobic enclosure and water-mediated effects \cite{CH7_Friesner2006}. In XP docking, the sampling algorithm is tightly coupled to XP scoring. Anchor fragments (often rings) from good SP poses are selected, and the ligand is regrown using an anchor-and-grow strategy. During regrowth, XP penalties and rewards guide focused sampling to avoid physically unrealistic poses and to exploit favorable hydrophobic pockets and hydrogen bond groups \cite{CH7_Friesner2006}. Because of this coupling, rescoring SP poses with XP GlideScore alone is generally not recommended; XP should be run as a full sampling + scoring protocol. While the full XP functional form is more complex, the top-level decomposition in XP score is written as \cite{CH7_Friesner2006}:

\[
  \text{XP\_GlideScore} \;=\; E_{\text{Coul}} \;+\; E_{\text{vdW}} \;+\; E_{\text{bind}} \;+\; E_{\text{penalty}}
\]

where:
\begin{align*}
  E_{\text{bind}} \;&=\; E_{\text{hyd} \cdot \text{enclosure}} \;+\; E_{\text{HBond} \cdot \text{nn} \cdot \text{motif}} \;+\; E_{\text{HBond} \cdot \text{cc} \cdot \text{motif}} \;+\; E_{\pi\pi} \;+\; E_{\text{HBond} \cdot \text{pair}} \;+\; E_{\text{phobic} \cdot \text{pair}} \\
  E_{\text{penalty}} \;&=\; E_{\text{desolv}} \;+\; E_{\text{buried\_polar}} \;+\; \cdots
\end{align*}

with:
\begin{itemize}
  \item $\text{E}_{\text{hyd} \cdot \text{enclosure}}$ is hydrophobic enclosure term, rewarding ligand groups that are well enclosed by lipophilic protein atoms on opposing faces.
  
  \item $\text{E}_{\text{HBond} \cdot \text{nn} \cdot \text{motif}}$, $\text{E}_{\text{HBond} \cdot \text{cc} \cdot \text{motif}}$ are specialized hydrogen bond motif terms for neutral–neutral and charged–charged interactions in hydrophobically enclosed environments respectively.
  
  \item $\text{E}_{\pi\pi}$ scores $\pi$–$\pi$ and other related aromatic stacking interactions involving the $\pi$ systems.
  
  \item $\text{E}_{\text{HBond} \cdot \text{pair}}$, $\text{E}_{\text{phobic} \cdot \text{pair}}$ scores the pairwise hydrogen bond and hydrophobic contacts respectively.
  
  \item $\text{E}_{\text{desolv}}$, $\text{E}_{\text{buried\_polar}}$, this includes explicit water-based solvation and buried polar penalties derived from docking water molecules into the complex and analyzing their interactions with polar/charged groups.

\end{itemize}


% edit from here
\section{AutoDock}
Autodock or commonly known as AD4, offers an easy to setup and run kind of environment for molecular docking experiments. Therefore, this becomes an easy tool for beginners wanting to learn molecular docking. The developers also separately provide all the tools/scripts, known as "MGLTools", necessary to set up systems for molecular docking. It is, therefore, an automated computational procedure primarily designed to predict the bound conformations of small, flexible ligands to macromolecular targets of known structure. Being an open-source tool with source code available, various groups have expanded the functionality of AD4 to improve docking accuracy for specific and challenging problems.\cite{CH7_Morris1998,CH7_Morris2009}. 
 It combines

\begin{itemize}
  \item a semi-empirical free-energy force field (the scoring function), and
  \item a global optimisation algorithm to search the space of ligand translations, orientations, and internal torsions. This enables the identification of low-energy binding poses. The original algorithm most commonly the Lamarckian genetic algorithm (LGA). However, later versions allower many conformational sampling methods.
\end{itemize}

AutoDock4 uses a grid-based strategy, as we have seen in the all the previously discussed method. The grid-based methods facilitates a significant speed up in for evaluating the receptor--ligand interaction energies by precomputing receptor contribution on three-dimensional grids (AutoGrid). During docking the energies are obtained by trilinear interpolation of these grids rather than explicit pairwise evaluation for each pose.\cite{CH7_Huey2007,CH7_Morris2009} This allows docking of many ligands or extensive conformational sampling of a single ligand at relatively modest computational cost. 

\subsection{The theory underlying AutoDock}

\subsubsection{The docking problem as a free-energy optimisation}

Conceptually, docking in AutoDock solves an optimisation problem:
\begin{quote}
  given a rigid (or partially flexible) receptor and a flexible ligand, find the ligand pose that minimises an empirical estimate of the binding free energy $\Delta G_{\text{binding}}$.
\end{quote}

A ligand pose is defined by
\begin{itemize}
  \item three Cartesian translation variables x, y, z
  \item three orientation variables $ \phi, \theta, \alpha $ (typically Euler angles or an equivalent rotation representation), and
  \item a set of torsional angles $ \psi_1, \dots, \psi_{N_{\text{tor}}}$ corresponding to the ligand's rotatable bonds.
\end{itemize}
Collectively, AutoDock represents these as a genotype vector
\[
  \Omega = \left\{ x, y, z, \phi, \theta, \alpha, \psi_1, \dots, \psi_{N_{\text{tor}}} \right\}.
\]
For any given \( \Omega \), AutoDock evaluates the scoring function \( f(\Omega) \), which approximates \( \Delta G_{\text{binding}} \). The task of docking is to find
\[
  \Omega^{\ast} = \arg\min_{\Omega} f(\Omega).
\]

\subsubsection{Representation of receptor and ligand}

In standard AutoDock 4:
\begin{itemize}
  \item The receptor is treated as rigid, although limited side-chain flexibility can be enabled by defining selected receptor residues as flexible (these flexible parts are technically treated as part of the ligand subsystem).\cite{CH7_Morris2009}
  \item The ligand is represented as a rigid root plus a tree of rotatable bonds; each rotatable bond defines a torsional degree of freedom $ \psi_k $.
  \item Coordinates, atom types and partial charges are stored in the PDBQT format, which extends PDB with AutoDock specific atom types and torsion tree information.
\end{itemize}

Atomic partial charges are typically Gasteiger charges, and atom types are assigned according to AutoDock's built-in atom typing scheme, which is essential for both the interaction grids and the desolvation model.\cite{CH7_Huey2007}

\subsubsection{Search algorithms: genetic algorithms and the Lamarckian GA}

AutoDock historically offered several global search methods (simulated annealing, genetic algorithms, and the Lamarckian genetic algorithm); for AutoDock 4, the default and recommended method is the Lamarckian genetic algorithm (LGA).\cite{CH7_Morris1998,CH7_Morris2009}

Key ideas of the LGA in AutoDock:
\begin{itemize}
  \item A population of candidate solutions (genotypes $ \Omega $) is maintained.
  \item Each candidate is mapped to a phenotype (3D ligand coordinates), scored by $ f(\Omega) $, and assigned a fitness (lower energy implies higher fitness).
  \item Standard genetic operators---selection, crossover, and mutation---are used to generate new candidate genotypes from fitter parents.
  \item A local search (e.g. Solis--Wets or a gradient-based method) is periodically applied to individual candidates to descend into nearby local minima of the scoring function.\cite{CH7_Morris1998,CH7_SantosMartins2021}
  \item Crucially, improvements found by the local search are written back into the genotype and thus inherited by descendants. This is the Lamarckian aspect: environmentally acquired traits (local minima) become heritable.
\end{itemize}

A typical LGA run:
\begin{enumerate}
  \item Initialise a population of random genotypes $ \Omega $.
  \item For each genotype in the population:
  \begin{enumerate}
    \item Decode to atomic coordinates (phenotype).
    \item Evaluate the scoring function $ f(\Omega) $.
  \end{enumerate}
  \item Apply selection: choose parent genotypes with probability proportional to fitness.
  \item Apply crossover and mutation to generate offspring genotypes.
  \item Apply local search to a fraction of the population; if local search improves a genotype, replace it with the improved genotype.
  \item Repeat evaluation, selection, crossover, mutation, and local search for a user\-defined number of generations or until convergence criteria are met.
  \item Return the best scoring individual from the run.
\end{enumerate}

In practice, a docking job consists of many independent LGA runs (e.g. 50), each starting from different random initial populations. The resulting best poses are then clustered by RMSD, and clusters ranked primarily by cluster size and cluster representative energy.\cite{CH7_Morris2009}

\subsection{The AutoDock4 Scoring Function}

\subsubsection{Semi-empirical free-energy force field}

AutoDock 4 uses a semi-empirical free-energy force field that estimates the binding free energy as a sum of physically motivated terms: van der Waals dispersion/repulsion, directional hydrogen bonding, electrostatics with a distance-dependent dielectric, a charge-based desolvation term, and an entropic penalty for ligand torsions.\cite{CH7_Huey2007,CH7_SantosMartins2021}

A convenient form of the AutoDock 4 scoring function is (units: kcal/mol)
\begin{equation}
\label{eq:ad4_scoring}
\begin{aligned}
  \Delta G_{\text{binding}} = {} &
  W_{\text{vdW}}
  \sum_{i,j}
  \left(
    \frac{A_{ij}}{s(r_{ij})^{12}}
    - \frac{B_{ij}}{s(r_{ij})^{6}}
  \right)
  \\
  & \;+\;
  W_{\text{hbond}}
  \sum_{i,j}
  E(t_{ij})
  \left(
    \frac{C_{ij}}{s(r_{ij})^{12}}
    - \frac{D_{ij}}{s(r_{ij})^{10}}
  \right)
  \\
  & \;+\;
  W_{\text{elec}}
  \sum_{i,j}
  \left(
    \frac{q_i \, q_j}{\varepsilon(r_{ij}) \, r_{ij}}
  \right)
  \\
  & \;+\;
  W_{\text{desolv}}
  \sum_{i,j}
  \left(
    \left( S_i V_j + S_j V_i \right)
    \exp\left[
      - \frac{r_{ij}^2}{2 \sigma^2}
    \right]
  \right)
  \\
  & \;+\;
  W_{\text{tor}} \, N_{\text{tor}}.
\end{aligned}
\end{equation}

Here:
\begin{itemize}
  \item  i and j  run over ligand and receptor atoms respectively (for intermolecular interactions), and over ligand atoms only for intramolecular ligand--ligand interactions.
  \item $ r_{ij} $ is the distance between atoms i  and  j .
  \item $ s(r_{ij}) $ is a smoothing function applied to soften the short-range repulsive wall of the 12--6 and 12--10 potentials (often implemented as a shifted distance, e.g. $ s(r_{ij}) = r_{ij} + \delta $, with $ \delta \approx 0.5 $ \r{A}).
  \item $ q_i$, $q_j$ are atomic partial charges (Gasteiger charges).
  \item $ \varepsilon(r_{ij}) $ is a distance-dependent dielectric function (commonly $ \varepsilon(r_{ij}) = \varepsilon_0 r_{ij} $, giving an effective screening that increases with distance).
  \item $S_i$ is the solvation parameter for atom  i, and $V_i$ is its atomic volume.
  \item $ \sigma$ is the width parameter of the Gaussian desolvation term (typically $ \sigma = 3.5$ \r{A}).
  \item $ N_{\text{tor}}$ is the number of rotatable bonds in the ligand.
  \item $ W_{\text{vdW}}$, $W_{\text{hbond}}$, $W_{\text{elec}}$, $W_{\text{desolv}}$, $W_{\text{tor}}$ are empirical weighting coefficients determined by linear regression against a training set of protein--ligand complexes with known binding affinities.
\end{itemize}

The energy is computed as the sum of intermolecular (ligand--receptor) and intramolecular (ligand--ligand) contributions; the latter enter via a precomputed estimate of the ligand's unbound state energy, so that the final reported score approximates the free energy difference between bound and unbound states.

\subsubsection{Interpretation of individual terms}
\begin{enumerate}
\item Van der Waals dispersion/repulsion

The first term,
\[
  \Delta G_{\text{vdW}} =
  W_{\text{vdW}}
  \sum_{i,j}
  \left(
    \frac{A_{ij}}{s(r_{ij})^{12}}
    - \frac{B_{ij}}{s(r_{ij})^{6}}
  \right),
\]
is a modified Lennard--Jones 12--6 potential. The coefficients $ A_{ij}$ and $B_{ij}$ depend on the atom types of i and j; they are chosen so that the minimum of the potential corresponds to the equilibrium contact distance and well depth appropriate to those atom types.\cite{CH7_Huey2007}

The $r^{-12}$ term models Pauli repulsion from overlapping electron clouds at very short distances, and the $r^{-6}$ term approximates London dispersion attraction. The smoothing function $s(r_{ij})$ broadens the attractive well and reduces numerical instability when atoms come too close.

\item Directional hydrogen bonding

The second term models hydrogen bonding interactions:
\[
  \Delta G_{\text{hbond}} =
  W_{\text{hbond}}
  \sum_{i,j}
  E(t_{ij})
  \left(
    \frac{C_{ij}}{s(r_{ij})^{12}}
    - \frac{D_{ij}}{s(r_{ij})^{10}}
  \right).
\]
Here $ C_{ij} $ and $ D_{ij}$ define a 12--10 potential specific to hydrogen bond donor--acceptor pairs, providing a narrower, more directional well than the generic van der Waals term.\cite{CH7_Huey2007} The angular term $E(t_{ij})$ depends on the hydrogen bond geometry (angle $ t_{ij} $); it down weights poorly aligned donor--acceptor pairs and favours near linear geometries.

\item Electrostatic interactions

The third term models Coulombic electrostatics:
\[
  \Delta G_{\text{elec}} =
  W_{\text{elec}}
  \sum_{i,j}
  \left(
    \frac{q_i \, q_j}{\varepsilon(r_{ij}) \, r_{ij}}
  \right).
\]
Partial charges $q_i$ are derived from the input structures and depend on the chosen charge model (commonly Gasteiger). The distance-dependent dielectric $ \varepsilon(r_{ij})$ mimics solvent screening without explicit solvent.\cite{CH7_Huey2007} Because this is a relatively simple electrostatic model, the associated weight  $W_{\text{elec}}$ is important to balance this term against the others when fitting to binding data.

\item Charge-based desolvation

The desolvation term accounts for the energetic cost of removing solvation from atoms upon binding, combined with favourable hydrophobic contacts:
\[
  \Delta G_{\text{desolv}} =
  W_{\text{desolv}}
  \sum_{i,j}
  \left(
    \left( S_i V_j + S_j V_i \right)
    \exp\left[
      - \frac{r_{ij}^2}{2 \sigma^2}
    \right]
  \right).
\]
The parameters $ S_i $ are solvation descriptors that depend on atom type and, via a charge\-based parametrisation, on the atomic charge.\cite{CH7_Huey2007} Positive $S_i$ correspond roughly to hydrophobic atoms (favourable desolvation on burial), whereas negative values correspond to polar/charged atoms (unfavourable desolvation).

The Gaussian factor localises the interaction to contact distances: for large $r_{ij}$, the exponential term is negligible (no desolvation interaction); for $ r_{ij} $ near contact distance, desolvation effects contribute significantly.

\item Torsional entropy penalty

Finally, the torsional entropy term
\[
  \Delta G_{\text{tor}} = W_{\text{tor}} \, N_{\text{tor}}
\]
penalises ligands with many rotatable bonds. When a ligand binds, many internal torsions become restricted, resulting in a loss of conformational entropy. AutoDock uses a simple linear model in the number of rotatable bonds $ N_{\text{tor}}$, with $W_{\text{tor}}$ obtained from regression against binding data.\cite{CH7_Huey2007}
\end{enumerate}

\subsubsection{Parameterisation and calibration}

The functional form of Eq.~\eqref{eq:ad4_scoring} was chosen based on physical intuition and computational efficiency, and the associated parameters $A_{ij}$, $B_{ij}$, $C_{ij}$, $D_{ij}$, $S_i$, $V_i$, $\sigma$ and weights $W_{\cdot}$ were fitted using linear regression to a dataset of protein--ligand complexes with known experimental binding free energies.\cite{CH7_Huey2007,CH7_Morris2009} The AutoDock 4 force field is thus semi-empirical: grounded in physical models but tuned to reproduce experimental data as closely as possible.

\subsection{A typical workflow for AutoDock4}

It is useful to divide the AutoDock workflow into four conceptual stages:

\begin{enumerate}
  \item \textbf{Structure preparation}
  \begin{enumerate}
    \item Receptor preparation
    \begin{enumerate}
  \item Obtain the receptor structure, typically from the Protein Data Bank (PDB).
  \item Clean the structure:
  \begin{itemize}
    \item Remove crystallographic water molecules unless specifically needed.
    \item Remove other non-relevant hetero atoms (e.g. buffer components, crystallisation agents).
    \item Check for missing residues or atoms and repair if necessary.
  \end{itemize}
  \item Add hydrogen atoms (particularly polar hydrogens) at a physically reasonable pH.
  \item Assign partial charges (e.g. Kollman or Gasteiger, depending on workflow/tool).
  \item Assign AutoDock atom types.
  \item Save the receptor in PDBQT format.
\end{enumerate}

This step is typically performed with AutoDockTools (ADT) or newer graphical front ends, but can also be scripted using e.g. MGLTools or external tools with PDBQT support.

  \item Ligand preparation
  \begin{enumerate}
  \item Start from a 2D or 3D structure of the ligand.
  \item Generate a reasonable 3D conformation (if starting from 2D), adding hydrogens and assigning bond orders.
  \item Assign protonation state and tautomers consistent with the target pH and binding site environment.
  \item Identify rotatable bonds and define the torsion tree:
  \begin{itemize}
    \item Choose a central root fragment (usually the largest rigid substructure).
    \item Define single bonds that are rotatable (excluding e.g. amide bonds, rings).
  \end{itemize}
  \item Assign partial charges (commonly Gasteiger charges in AutoDock workflows).
  \item Assign AutoDock ligand atom types.
  \item Save the ligand in PDBQT format.
\end{enumerate}
\end{enumerate}
  \item \textbf{Grid map generation (AutoGrid)}

AutoDock accelerates docking by precomputing three\-dimensional interaction energy maps for each ligand atom type using AutoGrid.\cite{CH7_Morris2009}

\begin{enumerate}
  \item Define the search space (grid box):
  \begin{itemize}
    \item Choose the grid center (often around a known binding site or co\-crystallised ligand).
    \item Choose grid dimensions (number of points in x, y, z and grid spacing, default $ \approx 0.375$ \r{A}).
    \item Ensure that the box fully encloses the binding site and allows some margin for ligand movement.
  \end{itemize}
  \item Generate a grid parameter file (GPF), specifying:
  \begin{itemize}
    \item Receptor PDBQT file.
    \item Grid box center and dimensions.
    \item List of ligand atom types for which to compute maps.
  \end{itemize}
  \item Run AutoGrid:
  \begin{itemize}
    \item For each ligand atom type, AutoGrid computes an interaction energy map over the 3D grid using the same functional forms as in Eq.~\eqref{eq:ad4_scoring}.
    \item Separate electrostatic and desolvation maps are also generated.
  \end{itemize}
\end{enumerate}

During docking, when the ligand is moved, interaction energies for ligand--receptor pairs are obtained by interpolating these maps at ligand atom positions, greatly reducing computational cost compared to evaluating all pairwise interactions explicitly for each pose.

  \item \textbf{Docking runs (AutoDock)}
  \begin{enumerate}
\item Docking parameter file (DPF)

Before running AutoDock, a docking parameter file (DPF) is prepared, specifying:
\begin{itemize}
  \item The ligand and receptor PDBQT files.
  \item The scoring function settings (typically defaults for AutoDock 4 force field).
  \item The search method (e.g LGA) and its parameters:
  \begin{itemize}
    \item Population size, number of generations.
    \item Number of energy evaluations.
    \item Mutation rate, crossover rate.
    \item Fraction of the population subjected to local search.
  \end{itemize}
  \item Number of independent docking runs (LGA runs).
\end{itemize}

\item Executing the docking

\begin{enumerate}
  \item Run AutoDock with the prepared DPF:
  \begin{itemize}
    \item AutoDock reads the receptor and ligand PDBQT files and the AutoGrid maps.
    \item It initialises random ligand poses and starts LGA runs according to the specified parameters.
  \end{itemize}
  \item AutoDock outputs a docking log file (DLG) containing:
  \begin{itemize}
    \item For each run: the best pose found, its estimated binding free energy $\Delta G_{\text{binding}}$, and the corresponding estimated inhibition constant $K_i$.
    \item Detailed information about the genetic algorithm progress (if enabled).
  \end{itemize}
\end{enumerate}
\end{enumerate}
  \item \textbf{Analysis and interpretation}
  After docking:
\begin{enumerate}
  \item Cluster the docked poses by RMSD (e.g. within 2 \r{A} threshold).
  \item Identify the largest clusters and their representative poses:
  \begin{itemize}
    \item Large, low-energy clusters suggest a well-defined binding mode.
    \item Small, scattered clusters may indicate multiple binding modes or insufficient sampling.
  \end{itemize}
  \item Inspect key interactions:
  \begin{itemize}
    \item Hydrogen bonds, salt bridges, hydrophobic contacts.
    \item Shape complementarity and burial of hydrophobic surfaces.
  \end{itemize}
  \item Use the predicted $\Delta G_{\text{binding}}$ values primarily for relative ranking of ligands rather than absolute binding affinity prediction; AutoDock's semi-empirical scoring tends to perform better for ranking within a series than for absolute affinities.\cite{CH7_Huey2007,CH7_Morris2009}
\end{enumerate}
\end{enumerate}



\section{AutoDock-Vina}

AutoDock Vina (usually referred to as Vina) is an open-source molecular docking program developed at The Scripps Research Institute to predict the binding mode and an approximate binding affinity of small molecules (ligands) to macromolecular receptors, typically proteins \cite{CH7_trott2010vina}.
Compared with its predecessor AutoDock 4, Vina introduces a new empirical scoring function, a continuous and differentiable potential tailored for gradient-based optimization, and efficient global search with multithreading support, leading to substantially improved speed and pose prediction accuracy \cite{CH7_trott2010vina,CH7_forli2016autodock}.
Because of its robustness and ease of use (automatic grid generation, relatively few parameters), Vina is widely used for training, small- to medium-scale docking, and virtual screening campaigns.

From a theoretical standpoint, docking can be seen as a compromise between physically detailed but expensive methods (e.g., molecular dynamics with explicit solvent) and fast but approximate approaches \cite{CH7_trott2010vina}.
Instead of explicitly simulating solvent and full receptor flexibility, Vina treats the receptor as largely rigid, the ligand as flexible around specified rotatable bonds, and uses a scoring function trained on experimental protein–ligand complexes (PDBbind) to approximate the standard chemical potentials (binding free energies) of complexes \cite{CH7_trott2010vina}.

\subsection{The theory underlying AutoDock-Vina}

\subsubsection{Conceptual view of molecular docking}

In molecular docking, the goal is to find ligand conformations and positions/orientations within a binding site that minimize an approximate binding free energy. 
Formally, one is interested in the minimum of a free-energy surface \( G(\mathbf{x}) \), where \( \mathbf{x} \) represents the collective coordinates of the ligand (and optionally some receptor degrees of freedom).
Exact evaluation of \( G(\mathbf{x}) \) requires integrating over solvent and many internal degrees of freedom, which is computationally prohibitive for routine screening \cite{CH7_sousa2006docking}.

Docking programs therefore introduce:
\begin{itemize}
  \item A \emph{reduced representation}: the receptor is mostly rigid; only ligand torsions (and sometimes selected side chains) are treated as flexible.
  \item An empirical \emph{scoring function}: a fast-to-evaluate function that correlates with experimental binding free energies.
  \item A stochastic \emph{search algorithm}: a global optimization method to explore the high-dimensional conformational space.
\end{itemize}

Vina follows this general framework but emphasizes (i) a smooth, differentiable scoring function that supports gradient-based local optimization; (ii) a compact set of empirically fitted parameters; and (iii) efficient use of multi-core CPUs via multithreading \cite{CH7_trott2010vina}.

\subsubsection{Representation of the docking problem in Vina}

In Vina, a docking pose (also called a conformation) is parameterized by:
\begin{itemize}
  \item The 3D translation of the ligand center.
  \item The 3D orientation of the ligand (represented internally, e.g., by a quaternion).
  \item The set of torsion angles around all active rotatable bonds in the ligand (and optionally in flexible side chains).
\end{itemize}

Given these degrees of freedom, the program:
\begin{enumerate}
  \item Builds the full 3D coordinates of all ligand atoms (and flexible atoms in the receptor).
  \item Evaluates an empirical scoring function c (described below) that depends on all pairwise interactions between atoms that can move relative to each other.
  \item Uses a global search algorithm (iterated local search) combined with a local optimizer (BFGS quasi-Newton) to drive the system towards low-scoring (low “energy”) conformations.
\end{enumerate}

The search is stochastic: multiple independent runs are started from random initial states, refined by local optimization, and then clustered and ranked according to the scoring function value.

\subsection{AutoDock Vina Scoring Function}

\subsubsection{General form}

The core of Vina is an empirical scoring function that is pairwise additive over atom pairs and designed to be continuous and differentiable almost everywhere. 
The conformation-dependent part of the scoring function is defined as \cite{CH7_trott2010vina}:
\[
  c = \sum_{i < j} f_{t_i t_j}\!\bigl(r_{ij}\bigr),
\]
where:
\begin{itemize}
  \item The sum is over all pairs of atoms \( i, j \) that can move with respect to each other, typically excluding 1–4 intramolecular pairs (atoms separated by three covalent bonds).
  \item \( t_i \) and \( t_j \) are atom types (following an X-Score-like typing scheme).
  \item \( r_{ij} \) is the interatomic distance between atoms \( i \) and \( j \).
  \item \( f_{t_i t_j}(r_{ij}) \) is an interaction function specific to the atom-type pair.
\end{itemize}

This total score \( c \) is the sum of intermolecular (\( c_{\text{inter}} \)) and intramolecular (\( c_{\text{intra}} \)) contributions:
\[
  c = c_{\text{inter}} + c_{\text{intra}}.
\]

During docking, Vina minimizes \( c \) (or, equivalently, searches for low values) to find favorable conformations. 
The final predicted binding affinity is then derived from the intermolecular component of the lowest-scoring pose and a simple entropic term penalizing ligand flexibility \cite{CH7_trott2010vina}.

\subsubsection{Surface distance and interaction functions}

Instead of expressing interactions directly as functions of \( r_{ij} \), Vina uses the surface distance between atoms:
\[
  d_{ij} = r_{ij} - R_{t_i} - R_{t_j},
\]
where \( R_{t_i} \) and \( R_{t_j} \) are van der Waals radii of atom types \( t_i \) and \( t_j \), respectively.

The pairwise interaction function is then written as
\[
  f_{t_i t_j}(r_{ij}) \equiv h_{t_i t_j}(d_{ij}),
\]
where \( h_{t_i t_j}(d) \) is a weighted sum of elementary terms:
\begin{itemize}
  \item Three generic steric terms (gaussian attractions and a quadratic repulsion), independent of atom types.
  \item A hydrophobic term, active only when both atoms are hydrophobic.
  \item A hydrogen-bond (H-bond) term, active only for donor–acceptor pairs.
\end{itemize}

All interaction functions are truncated beyond a cutoff distance (currently \( r_{ij} = 8 \,\text{\AA} \)) to limit computational cost \cite{CH7_trott2010vina}.

\subsubsection{Elementary interaction terms}

Vina defines the following steric terms as functions of \( d \) (in \AA) :
\[
  \text{gauss1}(d) = \exp\!\left( - \left( \frac{d}{0.5 \,\text{\AA}} \right)^{2} \right),
\]
\[
  \text{gauss2}(d) = \exp\!\left( - \left( \frac{d - 3 \,\text{\AA}}{2 \,\text{\AA}} \right)^{2} \right),
\]
\[
  \text{repulsion}(d) =
  \begin{cases}
    d^{2}, & d < 0, \\
    0,     & d \ge 0.
  \end{cases}
\]

The hydrophobic interaction term is a simple piecewise linear function of \( d \):
\[
  \text{hydrophobic}(d) =
  \begin{cases}
    1, & d < 0.5 \,\text{\AA}, \\
    \text{linear from 1 to 0}, & 0.5 \,\text{\AA} \le d \le 1.5 \,\text{\AA}, \\
    0, & d > 1.5 \,\text{\AA}.
  \end{cases}
\]

The hydrogen-bonding term is also piecewise linear, but active at shorter, slightly overlapping distances:
\[
  \text{hydrogen}(d) =
  \begin{cases}
    1, & d < -0.7 \,\text{\AA}, \\
    \text{linear from 1 to 0}, & -0.7 \,\text{\AA} \le d \le 0 \,\text{\AA}, \\
    0, & d > 0.
  \end{cases}
\]

Intuitively:
\begin{itemize}
  \item \(\text{gauss1}\) captures attractive contacts near van der Waals contact (\( d \approx 0 \)).
  \item \(\text{gauss2}\) captures medium-range attractive interactions (centered around \( d \approx 3 \,\text{\AA} \)).
  \item \(\text{repulsion}\) penalizes overlaps (negative \( d \), i.e., interpenetration).
  \item \(\text{hydrophobic}\) rewards burial of non-polar groups in close contact.
  \item \(\text{hydrogen}\) rewards favorable donor–acceptor separations typical of hydrogen bonds.
\end{itemize}

\subsubsection{Weights and final scoring expression}

The empirical scoring function combines these terms linearly. 
In the original Vina parameterization, the weights are \cite{CH7_trott2010vina}:
\[
  w_{\text{gauss1}} = -0.0356,\quad
  w_{\text{gauss2}} = -0.00516,\quad
  w_{\text{repulsion}} = 0.840,
\]
\[
  w_{\text{hydrophobic}} = -0.0351,\quad
  w_{\text{hydrogen}} = -0.587,
\]
and there is an additional weight associated with the number of rotatable bonds
\[
  w_{\text{rot}} = 0.0585.
\]

For a given conformation, the intermolecular contribution to the Vina score can be written schematically as:
\[
  c_{\text{inter}} =
    w_{\text{gauss1}} \sum_{i < j} \text{gauss1}\!\left(d_{ij}\right)
  + w_{\text{gauss2}} \sum_{i < j} \text{gauss2}\!\left(d_{ij}\right)
\]
\[
  +\, w_{\text{repulsion}} \sum_{i < j} \text{repulsion}\!\left(d_{ij}\right)
  + w_{\text{hydrophobic}} \sum_{i < j} \text{hydrophobic}\!\left(d_{ij}\right)
  + w_{\text{hydrogen}} \sum_{i < j} \text{hydrogen}\!\left(d_{ij}\right),
\]
where each sum is restricted to those atom pairs for which the corresponding term is defined (e.g., only hydrophobic atom pairs contribute to the hydrophobic term).

The intramolecular term \( c_{\text{intra}} \) has the same functional form but involves only interactions within the flexible parts of the ligand (and flexible receptor regions, if present). 
The conformation-dependent total score is
\[
  c = c_{\text{inter}} + c_{\text{intra}}.
\]

\subsubsection{Predicted binding affinity}

Vina converts the conformation-dependent score into a predicted binding affinity via a simple function g acting on the intermolecular part of the best-scoring pose and an entropic penalty for ligand flexibility \cite{CH7_trott2010vina}:
\[
  s_1 = g\!\left(c_{\text{inter}}^{(1)}\right) = c_{\text{inter}}^{(1)} + w_{\text{rot}} N_{\text{rot}},
\]
where:
\begin{itemize}
  \item \( c_{\text{inter}}^{(1)} \) is the intermolecular score of the lowest-\( c \) pose (pose 1).
  \item \( N_{\text{rot}} \) is the number of active rotatable bonds in the ligand.
  \item \( w_{\text{rot}} \) is the empirical weight given above.
\end{itemize}

The value \( s_1 \) is reported as the predicted binding free energy (in kcal/mol). 
It is important to emphasize that this is an empirical estimate: absolute values are approximate, and the scoring function is primarily optimized for pose prediction and ranking, not for quantitative free-energy prediction \cite{CH7_trott2010vina,CH7_forli2016autodock}.

For additional poses \( i \), Vina reports formal scores \( s_i \) that preserve the ranking while reusing the intramolecular term of the best pose:
\[
  s_i = g\!\left(c_i - c_{\text{intra}}^{(1)}\right).
\]

\subsection{Optimization Algorithm and Search Strategy}

Vina uses a hybrid global–local optimization procedure, combining:
\begin{itemize}
  \item \textbf{Iterated Local Search (ILS)} as the global search framework.
  \item \textbf{BFGS} (Broyden–Fletcher–Goldfarb–Shanno) as the local optimizer, which exploits gradients of the scoring function.
\end{itemize}

Each global-search “step” consists of:
\begin{enumerate}
  \item A random mutation of the current state (small change in translation, rotation, or torsion angles).
  \item Local minimization by BFGS, using both the score and its gradient.
  \item Acceptance or rejection according to a Metropolis-like criterion, allowing occasional uphill moves to escape local minima \cite{CH7_trott2010vina}.
\end{enumerate}

Multiple runs are performed from random starting conformations, and each run yields one or more low-scoring local minima. 
Vina then clusters these poses based on RMSD (typically using a cutoff such as 2 \r{A}) and reports cluster representatives ranked by predicted affinity.

Because the scoring function is smooth (continuous and differentiable, except for a few benign discontinuities), gradient-based methods are efficient and converge rapidly to local minima. 
Multithreading allows several independent search trajectories to run concurrently on different CPU cores \cite{CH7_trott2010vina}.

\subsection{Practical Workflow}

The following workflow summarizes how a typical user would perform a docking experiment with Vina, using either AutoDockTools (ADT)/ADFR tools or the newer Meeko-based scripts \cite{CH7_forli2016autodock,CH7_butt2020chimera}.

\subsubsection{Step 1: Receptor preparation}

\begin{enumerate}
  \item Obtain the 3D structure of the receptor from the Protein Data Bank (PDB) or homology modeling.
  \item Clean the structure:
  \begin{itemize}
    \item Remove crystallographic water molecules, ions, and cofactors not intended to be present in the docking calculation (or keep critical ones explicitly, e.g., metal ions).
    \item Check for missing residues or atoms; model them if necessary.
  \end{itemize}
  \item Add hydrogen atoms, especially polar hydrogens (important for hydrogen bonds and charge calculations).
  \item Assign atomic partial charges, e.g., Gasteiger charges (ADT/ADFR) or charges via Meeko-based workflows.
  \item Convert the receptor to PDBQT format, which encodes atom types, charges, and (for flexible residues) torsions.
\end{enumerate}

Modern workflows often use the Meeko scripts:
\begin{verbatim}
mk_prepare_receptor.py -i receptorH.pdb -o receptor -p -v \
    --box_size 20 20 20 --box_center x y z
\end{verbatim}
which generates a receptor PDBQT and a text file describing the docking box that can be used as a Vina configuration file.[cite:11]

\subsubsection{Step 2: Ligand preparation}

\begin{enumerate}
  \item Obtain or build a 3D structure of the ligand(s), preferably in a rich format (SDF, MOL2) that retains bond orders and formal charges.
  \item Determine the relevant protonation and tautomeric states (pH-dependent).
  \item Add hydrogen atoms and assign partial charges (e.g., Gasteiger).
  \item Define rotatable bonds (ADT or Meeko will do this automatically, subject to rules).
  \item Convert the ligand to PDBQT format:
\begin{verbatim}
mk_prepare_ligand.py -i ligand.sdf -o ligand.pdbqt
\end{verbatim}
  \item For multiple ligands (virtual screening), prepare each in PDBQT format with consistent chemistry and protonation conventions.
\end{enumerate}[cite:11]

\subsubsection{Step 3: Defining the search space (grid box)}

Vina requires a 3D search region (the “grid box”) within which the ligand is allowed to move. 
This is specified by:
\begin{itemize}
  \item Box center: \( \text{center\_x}, \text{center\_y}, \text{center\_z} \).
  \item Box dimensions: \( \text{size\_x}, \text{size\_y}, \text{size\_z} \) (in \AA).
\end{itemize}

Typical strategies:
\begin{itemize}
  \item If a co-crystallized ligand is available, choose a box centered on its binding site and extend the box margins (e.g., 10–15 \AA) to give the ligand room to explore.
  \item If the binding site is unknown, use pocket-detection tools or explore larger regions, at the cost of increased computation time.
\end{itemize}

A simple configuration file for Vina using the Vina force field might look like:
\begin{verbatim}
receptor = receptor.pdbqt
ligand   = ligand.pdbqt

center_x = 15.190
center_y = 53.903
center_z = 16.917

size_x   = 20.0
size_y   = 20.0
size_z   = 20.0

exhaustiveness = 8
num_modes      = 9
energy_range   = 3
\end{verbatim}[

\subsubsection{Step 4: Running AutoDock Vina}

With receptor and ligand PDBQT files and a configuration file in place, docking is initiated from the command line:
\begin{verbatim}
vina --config config.txt --out out_ligand_vina.pdbqt
\end{verbatim}

Key parameters:
\begin{itemize}
  \item \verb|--receptor| and \verb|--ligand|: override values in the config file if desired.
  \item \verb|--center_*| and \verb|--size_*|: specify the box directly on the command line.
  \item \verb|--exhaustiveness|: controls the thoroughness of the global search; default is 8, larger values such as 16–32 improve robustness at higher cost \cite{CH7_forli2016autodock}.
  \item \verb|--num_modes|: maximum number of output modes (poses) to report.
  \item \verb|--energy_range|: maximum energy difference (kcal/mol) between the best and worst reported poses.
\end{itemize}

The program prints a table of poses with:
\begin{itemize}
  \item Predicted affinity (kcal/mol).
  \item RMSD lower and upper bounds relative to the best pose.
\end{itemize}

\subsubsection{Step 5: Analyzing docking results}

\begin{enumerate}
  \item Load the receptor and output PDBQT into a molecular viewer (e.g., PyMOL, UCSF Chimera) or a docking-specific GUI (e.g., ADT, Chimera plugin).
  \item Inspect the top-ranked poses:
  \begin{itemize}
    \item Check that the ligand is in a chemically reasonable conformation.
    \item Verify key interactions (hydrogen bonds, salt bridges, hydrophobic contacts).
    \item Compare with known ligands or experimental SAR data when available.
  \end{itemize}
  \item For training and teaching, stepwise protocols such as those implemented in UCSF Chimera plus Vina provide an accessible workflow for beginners \cite{CH7_butt2020chimera}.
\end{enumerate}

\subsubsection{Interpretation and limitations}

Some important caveats and best practices:
\begin{itemize}
  \item The predicted binding energies are relative and approximate; they are most useful for ranking ligands against the same target, not as absolute free energies.
  \item Receptor rigidity is a strong approximation; induced fit and large-scale conformational changes are generally not captured.
  \item Docking results are sensitive to protonation states, tautomeric forms, and the definition of the search space.
  \item For critical projects, docking should be complemented by further refinement (e.g., MD simulations, MM/PBSA or MM/GBSA rescoring, or more rigorous free-energy calculations) \cite{CH7_forli2016autodock}.
\end{itemize}

Despite these limitations, AutoDock Vina provides an excellent platform for teaching the fundamentals of structure-based virtual screening, offering a clear conceptual link between molecular recognition, empirical scoring, and optimization algorithms.





%AutoDock Vina (Vina) employs a sophisticated method to sample ligand conformations, primarily through an efficient optimization algorithm that explores the ligand's positional, orientational, and torsional flexibility.
%Here's a breakdown of how AutoDock Vina samples ligand conformations:
%\begin{enumerate}
%\item Optimization Algorithm: Iterated Local Search (ILS)
%    \begin{enumerate}
%    \item Vina utilizes an Iterated Local Search global optimizer, which is similar to an Iterated Monte-Carlo (MC) search algorithm. This was chosen after exploring various stochastic global optimization approaches like genetic algorithms, particle swarm optimization, and simulated annealing.
%    \item The algorithm performs a succession of steps, each consisting of a mutation and a local optimization, with each step accepted according to the Metropolis criterion. This iterative process allows it to escape local minima and explore the conformational space more effectively.
%    \item Multiple runs are performed, starting from random conformations. These runs can be executed concurrently using multithreading on multicore CPUs, leveraging shared-memory hardware parallelism to speed up the process. The algorithm maintains a set of diverse significant minima found during these separate runs, which are then combined during the structure refinement and clustering stage. The number of steps in a run is determined adaptively based on the apparent complexity of the problem.
%    \end{enumerate}
%\item Local Optimization: BFGS Method with Gradients
%    \begin{enumerate}
%    \item For the local optimization steps within the Iterated Local Search, Vina uses the Broyden-Fletcher-Goldfarb-Shanno (BFGS) method. This is an efficient quasi-Newton method.
%    \item Crucially, BFGS uses not only the value of the scoring function but also its gradient. The gradient provides the derivatives of the scoring function for its arguments.
%    \item These arguments include the position and orientation of the ligand, as well as the values of the torsions for the active rotatable bonds in the ligand (and flexible residues, if any).
%    \item While evaluating the gradient in addition to the scoring function's value might take longer, using the gradient can significantly speed up the optimization process by giving the algorithm a "sense of direction".
%    \end{enumerate}
%\item Ligand Flexibility
%    Docking generally assumes that the receptor is rigid, while a chosen set of covalent bonds in the ligand are considered freely rotatable (referred to as active rotatable bonds). Vina can handle ligands with a varying number of active rotatable bonds, from 0 to 32, as demonstrated in tests. The number of heavy atoms and active rotatable bonds can influence the time taken for a Vina run.
%\item Scoring Function and Conformation Ranking
%    The optimization algorithm aims to find the global minimum of the scoring function's conformation-dependent part and other low-scoring conformations, which it then ranks. The scoring function itself, while not directly sampling conformations, guides the optimization process to identify favorable binding poses.
%\item Handling Search Space and Initial Conformations
%    \begin{enumerate}
%    \item Vina requires a user-defined "search space" in the receptor's coordinate system. For testing purposes, this space is typically expanded around the experimental ligand structure.
%    \item To prevent artificial help from a biased program, the initial ligand conformation (position, orientation, and torsions) is randomized to be unrelated to the experimental structure, while also avoiding self-clashes.
%    \end{enumerate}
%\item New Features in Vina 1.2.0 Enhancing Sampling
%    \begin{enumerate}
%    \item Macrocycle Conformational Sampling: Vina 1.2.0 introduced a specialized protocol for docking macrocycles, which are challenging due to correlated torsional changes. This method involves breaking one bond in the ring structure, using dummy atoms, and applying a linear potential to restore the bond. This allows independent exploration of torsional degrees of freedom and sampling of macrocycle conformations while adapting to the binding pocket, although at the cost of increased search complexity.
%    \item Control over Evaluations: Users can now specify the number of energy evaluations to be performed using the " $--\text{max}\_\text{evals}$ " option, providing more control over the search algorithm, particularly for complex searches (e.g., more flexible ligands or larger binding sites).
%    \end{enumerate}
%\end{enumerate}

%In summary, AutoDock Vina samples ligand conformations through an efficient Iterated Local Search, guided by a gradient-based BFGS local optimizer, which intelligently explores the ligand's translational, rotational, and torsional degrees of freedom. Recent updates further enhance its ability to handle complex flexibilities like those found in macrocycles.

%Molecular docking, including the method used by AutoDock Vina, generally operates under the assumption that much or all of the receptor is rigid. This approach represents a trade-off where representational detail is reduced in favor of computational speed.

%Here's how AutoDock Vina treats receptor flexibility:
%\begin{enumerate}
    
%\item Rigid Receptor Assumption: The foundational principle for Vina, as with many docking programs, is that the receptor structure remains constant during the docking process. For example, in tests evaluating Vina's performance, the receptors were consistently treated as rigid, while ligands were allowed to be flexible.
%\item Gradient-Based Optimization: While the primary focus of Vina's Broyden-Fletcher-Goldfarb-Shanno (BFGS) optimization method is on the ligand's position, orientation, and torsions, its arguments can also include the values of torsions for "flexible residues, if any". This indicates that Vina has a built-in capacity to account for some limited receptor flexibility (specifically, flexible residues) within its optimization algorithm, although it's not the default or primary mode of operation. The gradient calculation, which provides derivatives for these arguments, helps speed up the optimization process.
%\item Grid Maps and Post-processing: Vina automatically calculates its own grid maps based on the target (receptor) structure. These maps are used to define intermolecular interactions. Furthermore, after docking, Vina uses the receptor structure to perform a post-processing minimization of the docked poses. This refinement step relies on direct pairwise interactions with the receptor. The consistency of these operations implies a fixed receptor structure being used for these calculations.
%\item User Control and Default Behavior: Vina was designed with a limited amount of user input for performing docking, making it highly specialized and optimized. This design generally steers towards a rigid receptor assumption for simplicity and speed, as modifying receptor flexibility broadly would typically require significant changes in the source code.
%\item Comparison to AutoDock4: AutoDock4 (AD4) allowed users to modify a larger number of docking parameters, providing more direct access to engine internals, which made it suitable for developing new docking methods. However, even with AD4, the concept of full receptor flexibility was complex. AutoDock Vina 1.2.0, while maintaining Vina's core efficiency, introduced functionalities from the AutoDock Suite, some of which might implicitly or explicitly involve subtle receptor considerations (e.g., zinc coordination, which affects how the receptor interacts with specific ligand atoms), but the core assumption for bulk receptor movement generally remains rigid.
%\end{enumerate}

%The scoring function used in AutoDock Vina (referred to as Vina) is designed to approximate the standard chemical potentials of the system, aiming to reproduce chemical potentials rather than directly dealing with energies. This approach is more akin to "machine learning" than purely physics-based, being justified by its performance on test problems.
%The general functional form of the conformation-dependent part of the scoring function in AutoDock Vina is given by: 
%\begin{equation} \label{eq:vinascoring-1}
%\ \text{c} = \sum_{i<j} f_{t_i t_j} (r_{ij})
%\end{equation}
%• This summation is performed over all pairs of atoms that can move relative to each other, typically excluding 1–4 interactions (atoms separated by three consecutive covalent bonds).
%\begin{itemize}
%\item Each atom $i$ is assigned a type $t_i$.
%\item $f_{t_i t_j}$ represents a symmetric set of interaction functions based on the interatomic distance $r_{ij}$. Hydrogen atoms are not explicitly considered in this sum, other than for atom typing.
%This value, $c$, can be further broken down into intermolecular and intramolecular contributions: $\text{c} = \text{c}_{\text{inter}} + \text{c}_{\text{intra}}$
%The predicted free energy of binding ($s_1$) is calculated from the intermolecular part ($\text{c}_{\text{inter}1}$) of the lowest-scoring conformation: 
%\begin{equation} \label{eq:vinascoring-2}
%\begin{split}
%\ \text{s}_1 = \text{g}(\text{c}_1 - \text{c}_{\text{intra}1}) \\ 
%            = \text{g}(\text{c}_{\text{inter}1})
%\end{split}
%\end{equation}

%\item Here, $c_1$ refers to the total score of the lowest-scoring conformation, and $c_{\text{intra}1}$ is its intramolecular contribution.
%\item The function $g$ is an arbitrary, strictly increasing, smooth, and potentially nonlinear function.
%For other low-scoring conformations, $s_i$ values are formally given to preserve ranking, using the intramolecular contribution of the best binding mode ($\text{c}_{\text{intra}1}$): $\text{s}_i = \text{g}(\text{c}_i - c_{\text{intra}1})$
%\end{itemize}

%The specific conformation-independent function $g$ chosen for implementation is: 
%\begin{equation} \label{eq:vinascoring-3}
%\ \text{g}(\text{c}_{\text{inter}}) = \frac{\text{c}_{\text{inter}}}{1 + \text{wN}_{\text{rot}}}
%\end{equation}
%\begin{itemize}
%\item $N_{\text{rot}}$ represents the number of active rotatable bonds between heavy atoms in the ligand.
%\item $w$ is the associated weight for $N_{\text{rot}}$. As per Table 1, the weight for $N_{\text{rot}}$ is 0.0585.
%Components of the Interaction Functions ($f_{t_i t_j}$): The implementation of Vina's scoring function was largely inspired by X-score and tuned using the PDBbind dataset. The interaction functions $f_{t_i t_j}$ are defined relative to the surface distance $d_{ij} = r_{ij} - R_{t_i} - R_{t_j}$, where $R_t$ is the van der Waals radius of atom type $t$. In Vina's scoring function, $h_{t_i t_j}$ (which $f_{t_i t_j}$ is equivalent to) is a weighted sum of steric interactions, hydrophobic interactions, and hydrogen bonding.
%The specific terms and their weights used in Vina's scoring function are provided in Table 1:
%\item Steric Terms: These are identical for all atom pairs.
%\begin{itemize}
%    \item gauss1: weighted by -0.0356. $\text{gauss1}(\text{d}) = e^{-(\text{d}/0.5\text{ Å})^2}$
%    \item gauss2: weighted by -0.00516. $\text{gauss2}(\text{d}) = e^{-((\text{d}-3\text{ Å})/2\text{ Å})^2}$
%    \item Repulsion: weighted by 0.840. $ \text{repulsion}(\text{d}) = 
%                        \begin{cases} 
%                        \text{d}^2, & \text{if } \text{d} < 0 \cr 0, & \text{if } \text{d} \geq 0 
%                        \end{cases}$
%\end{itemize}
%\item Hydrophobic Term: weighted by -0.0351. It equals 1 when $\text{d} < 0.5${\AA}, 0 when $\text{d} > 1.5${\AA}, and is linearly interpolated between these values.
%\item Hydrogen Bonding Term: weighted by -0.587. It equals 1 when $\text{d} < -0.7${\AA}, 0 when $\text{d} > 0$, and is similarly linearly interpolated in between. Metals are formally treated as hydrogen bond donors.
%\item All interaction functions $\text{f}_{\text{t}_i \text{t}_j}$ are cut off at $\text{r}_{ij} = 8${\AA}
%\end{itemize}

%Comparison with AutoDock4 (AD4) Scoring Function The Vina scoring function is quite different from the AD4 scoring function.
%\begin{enumerate}
%\item AD4 scoring function uses a physics-based model that includes van der Waals, electrostatic, and directional hydrogen-bond potentials derived from an early AMBER force field, along with a pairwise-additive desolvation term based on partial charges, and a simple conformational entropy penalty.
%\item In contrast, Vina's scoring function notably lacks explicit electrostatics and solvation terms. It relies on a van der Waals-like potential (combining repulsion and two attractive Gaussians), a non-directional hydrogen-bond term, a hydrophobic term, and a conformational entropy penalty.
%\item The computational cost for energy evaluations is higher with the AD4 scoring function in Vina 1.2.0, taking nearly three times longer than with the Vina scoring function, due to the need to interpolate additional electrostatic and desolvation maps for each movable atom. Vina ranks conformations based on its scoring function, which includes both intermolecular and intramolecular contributions, differing from X-score, which only considers intermolecular contributions.
%\end{enumerate}

\section{FlexX}
FLEXX performs ligand conformational search for docking by modelling ligand flexibility using a discrete model. This approach focuses on sampling the conformational space of the ligand efficiently.
Here's a breakdown of how it handles ligand flexibility:
\begin{enumerate}
\item Fixed Geometries: Bond lengths and angles are kept invariant as given in the input structure, except within ring systems. Therefore, the ligand structures provided to FLEXX should ideally be reasonably minimised geometries.
\item Acyclic Single Bond Flexibility (Torsion Angles):
    \begin{enumerate}
    \item For acyclic single bonds, a set of preferred torsion angles is assigned.
    \item These preferred angles are determined by a statistical evaluation of the Cambridge Structural Database (CSD).
    \item The method uses a database of approximately 900 molecular fragments with a central single bond. By matching these representative torsional fragments onto the ligand, a histogram of torsion angle occurrences from the CSD is used to select highly populated torsion angles for generating low-energy conformations.
    \item Up to 12 discrete torsion angle values can be assigned to each acyclic single bond.
    \item Bonds adjacent to a methyl group or a planar amino group are exceptions and remain unchanged from the initial conformation.
    \end{enumerate}
\item Ring Flexibility:
    \begin{enumerate}
    \item Multiple conformations for ring systems are computed using external programs.
    \item Initially, the SCA program (developed by De Clercq) was used, which computes alternative conformations for polycyclic systems containing five- to seven-membered rings. This was considered sufficient for most practical purposes in drug design and typically took only a few seconds for a typical drug.
    \item Later versions of FLEXX (e.g., version 1.6.5 evaluated in Kramer et al., 1999) used the CORINA program for computing multiple ring conformations, with a limit of seven ring atoms. Larger rings were treated as rigid based on the input structure.
    \end{enumerate}
\item Integration with Docking Algorithm: This discrete conformational model is fundamental to FLEXX's incremental construction algorithm. The ligand is fragmented, and then fragments are added step by step in all possible conformations, allowing the system to explore the ligand's flexible states as it builds the complex within the active site.
\item Accuracy and Complexity: The described model of ligand flexibility is reported to cause RMS deviations only in terms of conformational differences of less than 1 \r{A} (e.g., 0.4 \r{A} for methotrexate). While this model generates a combinatorial number of conformations, internal clashes discard a portion of these, providing an upper bound for the number of possible low-energy conformations. The ability to handle ligands with up to 17 rotatable bonds is sufficient for most practical drug design purposes. However, for ligands with more than 20 fragments, other search strategies might be required due to the heuristic greedy strategy employed during complex construction.
\item Pre-processing: Before docking, ligand structures are energy-minimised and assigned correct atom types, bond types, and formal charges to ensure suitable bond distances and angles and to remove any "docking information" from the PDB structure.
\end{enumerate}
This systematic, discrete approach to ligand flexibility allows FlexX to efficiently sample a large variety of ligand conformations, aiming to reproduce the experimentally observed binding mode.

\subsection{Scoring function}
FLEXX scores the docked ligands using an empirical scoring function similar to that developed by Böhm (1994), with some minor changes. This function is designed to estimate the free binding energy (DG) of the protein-ligand complex.
The general equation for the scoring function is:
The equation for the scoring function is:
\begin{equation} \label{eq:flexscoring}
\begin{split}
\ \Delta \text{G} = \\ 
                    \Delta \text{G}_0 + \Delta\text{G}_{\text{rot}} \times \text{N}_{\text{rot}} + \\
                    \Delta\text{G}_{\text{hb}}\sum_{\text{neutral H-bonds}} f(D_R, D_\alpha) + \\
                    \Delta\text{G}_{\text{io}}\sum_{\text{ionic int.}} f(D_R, D_\alpha) + \\ 
                    \Delta \text{G}_{\text{aro}}\sum_{\text{aro int.}} f(D_R, D_\alpha) + \\
                    \Delta \text{G}_{\text{lipo}}\sum_{\text{lipo. cont.}} f^*(D_R) 
\end{split}
\end{equation}
Here's a breakdown of each term:
\begin{itemize}
\item \textbf{$\Delta\text{G}_0$}: This is a fixed ground term, set at 5.4 kJ/mol.

\item \textbf{$\Delta G_{\text{rot}} \times N_{\text{rot}}$}: $\text{N}_{\text{rot}}$ represents the number of free rotatable bonds in the ligand that become immobilized (hindered) upon binding to the complex. $\Delta \text{G}_{\text{rot}}$ accounts for the loss of entropy during ligand binding due to this immobilisation and is valued at 1.4 kJ/mol.

\item \textbf{$\Delta \text{G}_{\text{hb}} \sum_{\text{neutral H-bonds}} f(D_R, D_\alpha)$}: This term is a sum over all pairwise neutral hydrogen bonds. $f(D_R, D_\alpha)$ is a scaling function that penalizes deviations from ideal geometry for these interactions. The ideal interaction distance for hydrogen bonds is 1.9 \r{A}.

\item \textbf{$\Delta \text{G}_{\text{io}} \sum_{\text{ionic int.}} f(D_R, D_\alpha)$}: This sums over all pairwise ionic interactions. $f(D_R, D_\alpha)$ is the same scaling function for geometric deviations. The ideal interaction distance for metal interactions is 2.0 \r{A}.

\item \textbf{$\Delta \text{G}_{\text{aro}} \sum_{\text{aro int.}} f(D_R, D_\alpha)$}: This term sums over aromatic interactions, specifically involving interactions between aromatic-ring-atoms, methyl, or amide groups with an aromatic-ring-centre. These interactions are used to generate the preferred t-shaped arrangement of neighbouring aromatic ring systems. In some versions, an additional parameter $\Delta \text{G}_{\text{aro}}$ of -0.7 kJ/mol is included. The ideal interaction distance for aromatic interactions is 4.5 \r{A}.

\item \textbf{$\Delta \text{G}_{\text{lipo}} \sum_{\text{lipo. cont.}} f^*(D_R)$}: This is a modification of Böhm’s original lipophilic contact energy, which was estimated using a grid method. In FLEXX, it's calculated as a sum over all pairwise atom-atom contacts.  \text{$f^*(D_R)$} is a function that accounts for contacts with ideal distances and penalises forbiddingly close contacts. \text{$D_R = R - R_0$}, where R is the actual distance between atom centers and \text{$R_0$} is its ideal value (sum of van-der-Waals radii, each increased by 0.3 \r{A}). The function $f^*(D_R)$ has specific ranges for its values.

\end{itemize}

The model of protein-ligand interactions and the scoring function in FLEXX are derived from the work of Böhm (1992a, 1994) and Klebe (1994). The scoring function is computationally fast and easy to compute, contributing to FLEXX's short runtimes. Its primary purpose is for ranking the generated docking solutions.
While useful for distinguishing between good and bad inhibitors, the absolute binding energy values often show large differences (about 20 kJ/mol) from experimental values. This indicates that some important energetic contributions may not be fully modeled, and improvements to the scoring function are an ongoing area of development. For instance, it may not penalize large hydrophobic groups exposed to the solvent, potentially leading to misranking. Filter functions, which discard placements with undesirable properties like insufficient steric complementarity or exposed lipophilic regions, are being developed to address these issues.

Key Aspects of Scoring in FLEXX:
\begin{itemize}
\item Ranking Solutions: The scoring function is primarily used for ranking the generated docking solutions.
\item Origin of Chemical Model: The underlying chemical model in FLEXX is based on the work of Böhm (1992a, 1994) and Klebe \& Mietzner (1994).
\item Interaction Types: Interactions are classified by their geometric constraints and are used for placing the ligand. These include hydrogen bonds and some hydrophobic interactions (e.g., phenyl rings and methyl groups). These interactions are arranged on three levels, from highly directional (Level 3, e.g., H-bonds, metal acceptors) down to directionally unspecific (Level 1, e.g., aliphatic and aromatic carbon atoms, sulfur).
\item Heuristic Nature: The functions f and f* are described as heuristic distance and angle-dependent penalties.
\item Scoring Speed: The scoring function is chosen to be fast and easy to compute, contributing to FLEXX's extremely short runtimes. Docking a typical drug molecule, considering ligand flexibility, can be performed in two to three minutes on a standard workstation, making it approximately ten times faster than alternative approaches for flexible ligand docking at the time of its development.
\item Limitations and Improvements:
\begin{itemize}
    \item While useful for roughly distinguishing good and bad inhibitors, the absolute binding energy values often show large differences (about 20 kJ/mol) from experimental values.
    \item The scoring function may not penalise large hydrophobic groups reaching out into the solvent, potentially leading to misranking of solutions.
    \item Filter functions that extract placements with "undesirable properties" (e.g., insufficient steric complementarity or large lipophilic regions exposed to solvent) have been developed to improve ranking and quality.
    \item There is ongoing work to improve the scoring function to better model important energetic contributions that are currently not accounted for. Using a simple function based on a single configuration of a receptor-ligand complex for scoring is acknowledged as a rough approximation of the free energy of binding.
    \item Including water molecules explicitly is crucial for reproducing correct binding modes in some cases, which influences the interaction pattern and thus the score.
\end{itemize}
\end{itemize}

\section{SurFlex}

Surflex-Dock is a widely used small-molecule docking program that combines an empirically parameterized scoring function with a surface-based similarity search strategy to generate and rank protein--ligand binding modes.\cite{CH7_Jain2003Surflex,CH7_Jain2007Surflex21} It was originally introduced as Surflex,\cite{CH7_Jain2003Surflex} implemented in the SYBYL environment, and has since been extended to Surflex-Dock 2.1 and beyond.\cite{CH7_Jain2007Surflex21}

Two key design principles underlie Surflex-Dock:
\begin{itemize}
  \item An empirical scoring function that outputs values in units of $pK_d$, trained on experimental protein--ligand complexes and designed to be smooth, continuous, and piecewise differentiable for efficient gradient-based optimization.\cite{CH7_Jain2007Surflex21,CH7_Jain2006NegTrain}
  \item A molecular similarity-based search centered on an idealized protomol representation of the binding site, which allows rapid generation of candidate poses by aligning ligand fragments to this protomol.\cite{CH7_Jain2003Surflex}
\end{itemize}

This document gives a conceptual overview of Surflex-Dock, its typical workflow, the theoretical basis of its scoring and search strategies, and a detailed explanation of the scoring function with equations suitable for training beginning users.

\subsection{Search and optimization}

Surflex-Dock 2.1 introduced several important enhancements to the search protocol:\cite{CH7_Jain2007Surflex21}

\begin{itemize}
  \item \textbf{Quasi-Newton optimization}: Local pose optimization uses the BFGS quasi-Newton method with analytical gradients of the scoring function, improving both speed and robustness over simple steepest-descent approaches.
  \item \textbf{Small-molecule force field}: An implementation of the DREIDING force field is used to model intramolecular ligand energetics (bond lengths, angles, torsions, van der Waals interactions). This allows optimization of Cartesian coordinates beyond pure pose-parameter (translation, rotation, torsion) space.
  \item \textbf{Ring flexibility}: General dynamic exploration of flexible ring conformations (five-, six-, and seven-membered rings) is supported by mapping prototype ring conformers (e.g., cyclohexane) onto ligand ring templates and minimizing with the force field.
  \item \textbf{Fragment-guided docking}: When a particular substructure is known to bind in a defined manner (e.g., hinge-binding motifs in kinases, chelating groups in metalloenzymes), this fragment can be pre-placed and used to constrain and accelerate pose generation for analogs.
\end{itemize}

These extensions improve sampling of relevant poses, particularly for ligands with significant ring flexibility or initial conformational strain, and enhance screening robustness across diverse ligand sets.\cite{CH7_Jain2007Surflex21}


\subsection{Scoring Function}

Surflex-Dock uses an empirical scoring function that is constructed as a sum of terms corresponding to physically motivated contributions to binding:
\begin{itemize}
  \item hydrophobic (steric) complementarity,
  \item polar (hydrogen-bond and electrostatic) complementarity,
  \item entropic penalties for fixing ligand degrees of freedom,
  \item solvation-related corrections.
\end{itemize}

Instead of computing an absolute binding free energy from first principles, the function is parameterized against a training set of protein--ligand complexes of known affinity and structure. The output is directly in units of
\[
  pK_d = -\log_{10}(K_d),
\]
so that a one-unit change in the Surflex score roughly corresponds to an order-of-magnitude change in dissociation constant within the regime of the training data.\cite{CH7_Jain2007Surflex21,CH7_Jain2006NegTrain}

A notable methodological feature is the use of a multiple-instance learning framework during parameter estimation.\cite{CH7_Jain2007Surflex21} Rather than assuming that the exact crystallographic pose defines the score, the method:
\begin{enumerate}
  \item Identifies a local optimum of the scoring function in the vicinity of the crystallographic pose.
  \item Associates the complex's experimental affinity with this locally optimal pose.
\end{enumerate}
This reflects the idea that relatively small perturbations of a crystallographic pose are often experimentally indistinguishable, but can have large nominal effects on a given scoring function.

Subsequent work introduced the use of synthetically generated negative training data (physically unreasonable, interpenetrating poses) so that parameters controlling clash penalties and other ``bad'' interactions are estimated rigorously rather than heuristically.\cite{CH7_Jain2006NegTrain}

\subsection{Protomol and molecular similarity-based search}

A central concept in Surflex-Dock is the protomol.\cite{CH7_Jain2003Surflex}
\begin{itemize}
  \item The protomol is an idealized ligand-like object that occupies the binding site.
  \item It is constructed from small fragments such as methane (CH\(_4\)), amide carbonyl (C=O), and amide N--H, tessellated over the binding pocket and optimized using the Surflex scoring function itself.
  \item The protomol thus encodes the locations of preferred hydrophobic contacts, hydrogen bond donors, and hydrogen bond acceptors, as well as steric boundaries.
\end{itemize}

Docking a new ligand proceeds via surface-based molecular similarity:
\begin{enumerate}
  \item Fragments of the input ligand are generated (e.g., by breaking rotatable bonds).
  \item These fragments are aligned to the protomol using a rapid morphological similarity measure that compares molecular surfaces and chemical features.
  \item High-similarity fragment poses are retained and locally optimized using the Surflex scoring function.
  \item Compatible, high-scoring fragments are merged to generate full-ligand poses, which are further refined by gradient-based optimization.
\end{enumerate}

Compared to purely grid-based or exhaustive sampling approaches, this strategy greatly reduces the search space by focusing on poses that resemble an idealized ligand for the pocket.


\subsection{Typical Surflex-Dock Workflow}

\begin{enumerate}

\item Step 1: Receptor preparation

\begin{enumerate}
  \item \textbf{Select and clean the protein structure}
  \begin{itemize}
    \item Choose an appropriate PDB structure (good resolution, relevant conformation, appropriate cofactors and metal ions).
    \item Remove crystallographic artifacts (e.g., irrelevant buffer molecules) while retaining essential structural waters, ions, and cofactors as needed.
  \end{itemize}
  \item \textbf{Add hydrogens and assign protonation states}
  \begin{itemize}
    \item Add all hydrogens, not just polar hydrogens.
    \item Assign protonation and tautomeric states for ionizable residues and active-site residues based on pH and the known mechanism.
  \end{itemize}
  \item \textbf{Define the binding site}
  \begin{itemize}
    \item If a co-crystallized ligand is present, define the binding site as a region around that ligand (e.g., within a chosen radius).
    \item Alternatively, define the site with a binding pocket-finding algorithm or by specifying residues.
  \end{itemize}
\end{enumerate}

\item Step 2: Protomol generation

Surflex can generate the protomol in two main ways:\cite{CH7_Jain2003Surflex}
\begin{itemize}
  \item \textbf{Ligand-based protomol}: When a co-crystallized ligand is available, the protomol is grown to mimic its interactions, essentially capturing an idealized version of that ligand's binding mode.
  \item \textbf{Pocket-based protomol}: When no ligand is available, the protomol can be constructed directly from the pocket geometry using fragment tessellation and optimization.
\end{itemize}

For training, a ligand-based protomol is often recommended because it usually yields more intuitive results and is easier to visualize.

\item Step 3: Ligand preparation

\begin{enumerate}
  \item Generate 2D or 3D structures of ligands (from SMILES, SD files, etc.).
  \item Standardize valence, formal charges, and tautomers as appropriate.
  \item Assign protonation states consistent with the docking pH.
  \item Optionally, pre-minimize ligands using the built-in DREIDING implementation to reduce conformational strain before docking.
\end{enumerate}

\item Step 4: Docking setup and execution

Key Surflex-Dock options for beginners typically include:
\begin{itemize}
  \item \textbf{Sampling level} (e.g., standard vs.\ thorough sampling).
  \item \textbf{Ring flexibility} (enable the ring-search option for ligands with flexible ring systems).
  \item \textbf{Fragment-guided docking} when a known binding substructure is available (using the appropriate fragment-matching options).
  \item \textbf{Number of poses per ligand} to be retained (e.g., top 10--20 poses).
\end{itemize}

The docking engine:
\begin{enumerate}
  \item Aligns ligand fragments to the protomol using surface similarity.
  \item Applies thresholds to eliminate gross steric clashes and low-similarity poses.
  \item Merges compatible high-scoring fragments into full ligand poses.
  \item Optimizes each pose using BFGS with respect to the Surflex scoring function.
  \item Returns a ranked list of poses with associated Surflex scores (in units of \(pK_d\)).
\end{enumerate}

\item Step 5: Post-processing and analysis

After docking:
\begin{itemize}
  \item Inspect top-ranked poses visually for key interactions (hydrogen bonds, hydrophobic contacts, metal coordination).
  \item Cluster poses by root-mean-square deviation (RMSD) to identify distinct binding modes.
  \item Consider rescoring with alternative scoring functions or consensus scoring when prioritizing compounds for experimental testing.
\end{itemize}
\end{enumerate}

\subsection{The Surflex-Dock Scoring Function}

\subsubsection{Atomic surface distance}

The Surflex-Dock scoring function is formulated in terms of atomic surface distances between protein and ligand atoms.\cite{CH7_Jain2007Surflex21}

For a ligand atom i and a protein atom j, define:
\[
  r_{ij} = d_{ij} - (R_i + R_j),
\]
where:
\begin{itemize}
  \item \(d_{ij}\) is the center-to-center distance between atoms i and j,
  \item $R_i$ and $R_j$ are their van der Waals radii.
\end{itemize}

Thus:
\begin{itemize}
  \item \(r_{ij} > 0\) corresponds to separated atoms (no overlap),
  \item \(r_{ij} = 0\) corresponds to exact contact of van der Waals surfaces,
  \item \(r_{ij} < 0\) corresponds to interpenetration (steric clash or close contact).
\end{itemize}

Atoms are typed as \emph{nonpolar} (e.g., carbon and attached hydrogens) or \emph{polar} (e.g., heteroatoms involved in hydrogen bonding). Polar atoms may also carry formal charges.

\subsubsection{Hydrophobic (steric) complementarity term}

For each relevant atom pair, Surflex defines a steric (hydrophobic) complementarity score as a function of the surface distance r (writing r instead of $r_{ij}$) for brevity):\cite{CH7_Jain2007Surflex21}
\begin{equation}
  \text{steric\_score}(r)
  =
  \lambda_{1} \exp\!\left(
    -\frac{\left(r + n_{1}\right)^{2}}{n_{2}}
  \right)
  +
  \frac{\lambda_{2}}{
    1 + \exp\!\left(
      n_{3} \left( r + n_{4} \right)
    \right)
  }
  +
  \lambda_{3} \max\!\left( 0,\left( r + n_{5} \right)^{2} \right).
  \label{eq:steric}
\end{equation}

Here:
\begin{itemize}
  \item \(\lambda_{1}, \lambda_{2}, \lambda_{3}\) and \(n_{1}, \dots, n_{5}\) are empirically fitted parameters.
  \item The first term is a Gaussian centered near the optimal contact distance between hydrophobic atoms.
  \item The second term is a sigmoidal function controlling the mid-range attractive behavior.
  \item The third term is a quadratic penalty that activates for sufficiently positive r (i.e., when atoms are too close or clashing, depending on the shift $n_5$).
\end{itemize}

Physically, for ideal hydrophobic contacts, \(\text{steric\_score}(r)\) contributes on the order of \(0.1\) units of $pK_d$ per favorable nonpolar atom--atom contact, with rapidly decaying contributions as the atoms move apart or interpenetrate excessively.\cite{CH7_Jain2007Surflex21}

\subsubsection{Polar complementarity term}

The polar complementarity term has a similar functional form but is applied to polar atom pairs and includes additional scaling to account for charge and directionality:\cite{CH7_Jain2007Surflex21}
\begin{equation}
  \text{polar\_score}(r)
  =
  \lambda_{4} \exp\!\left(
    -\frac{\left(r + n_{6}\right)^{2}}{n_{7}}
  \right)
  +
  \frac{\lambda_{5}}{
    1 + \exp\!\left(
      n_{3} \left( r + n_{8} \right)
    \right)
  }
  +
  \lambda_{3} \max\!\left( 0,\left( r + n_{9} \right)^{2} \right).
  \label{eq:polar}
\end{equation}

Key aspects:
\begin{itemize}
  \item \(\lambda_{4}, \lambda_{5}\) and \(n_{6}, \dots, n_{9}\) are also empirically fitted.
  \item The parameter \(\lambda_{3}\) is shared as the coefficient for the quadratic clash term in both hydrophobic and polar interactions, reflecting a common penalty model for interpenetration.
  \item The polar term is further scaled by:
  \begin{itemize}
    \item a directionality factor that depends on the relative geometry (e.g., donor--hydrogen--acceptor angle and orientation),
    \item a factor that depends quadratically on the formal charges of the interacting atoms.
  \end{itemize}
\end{itemize}

For an ideal hydrogen bond (e.g., between a donor proton and a carbonyl oxygen), this term can contribute on the order of \(1.2\) units of \(pK_d\), with an optimal surface separation corresponding to a typical donor--acceptor distance near \(1.97\) \AA.\cite{CH7_Jain2007Surflex21}

\subsubsection{Total Surflex-Dock score}

The overall Surflex-Dock score for a given protein--ligand pose is a sum over pairwise atomic terms plus global entropic and solvation contributions:\cite{CH7_Jain2007Surflex21,CH7_Jain2006NegTrain}
\begin{equation}
  S_{\text{Surflex}}
  =
  \sum_{(i,j) \in \mathcal{N}_{\text{np}}}
    \text{steric\_score}\!\left( r_{ij} \right)
  +
  \sum_{(i,j) \in \mathcal{N}_{\text{pol}}}
    w_{ij}^{\text{pol}} \,
    \text{polar\_score}\!\left( r_{ij} \right)
  +
  S_{\text{ent}}
  +
  S_{\text{solv}}.
  \label{eq:surflex_total}
\end{equation}

In this expression:
\begin{itemize}
  \item \(\mathcal{N}_{\text{np}}\) is the set of nonpolar protein--ligand atom pairs that are considered for steric (hydrophobic) scoring.
  \item \(\mathcal{N}_{\text{pol}}\) is the set of polar protein--ligand atom pairs.
  \item \(w_{ij}^{\text{pol}}\) is a weighting factor that encodes:
  \begin{itemize}
    \item directional dependence of the polar interaction (e.g., hydrogen-bond angle),
    \item scaling with formal charges of atoms \(i\) and \(j\).
  \end{itemize}
  \item \(S_{\text{ent}}\) represents entropic penalties:
  \begin{equation}
    S_{\text{ent}}
    =
    c_{\text{rot}} N_{\text{rot}}
    +
    c_{\text{tor}} N_{\text{tor}}
    +
    c_{\text{tr}} \log\!\left(\text{MW}\right),
    \label{eq:entropic}
  \end{equation}
  where:
  \begin{itemize}
    \item \(N_{\text{rot}}\) is the number of rotatable bonds in the ligand,
    \item \(N_{\text{tor}}\) may count additional torsional contributions,
    \item \(\text{MW}\) is the molecular weight of the ligand,
    \item \(c_{\text{rot}}, c_{\text{tor}}, c_{\text{tr}}\) are empirically estimated coefficients.
  \end{itemize}
  \item \(S_{\text{solv}}\) is a solvation-related correction term, which in the original Surflex-Dock parameterization has relatively small weight but can be refined or customized.\cite{CH7_Jain2007Surflex21}
\end{itemize}

All coefficients and parameters (\(\lambda_{k}\), \(n_{k}\), and \(c_{k}\)) are obtained by fitting to training sets of protein--ligand complexes with known affinities using the multiple-instance learning protocol. Later refinements incorporate negative training data (unphysical poses) to better constrain the parameters governing interpenetration and other unfavorable interactions.\cite{CH7_Jain2006NegTrain}

\subsection{Interpretation of Surflex scores}

Because \(S_{\text{Surflex}}\) is in units of \(pK_d\),
\begin{itemize}
  \item A higher score indicates stronger predicted binding affinity.
  \item Differences in score can be roughly interpreted as differences in log affinity, e.g., a difference of approximately 1 unit in score corresponds to roughly a tenfold change in predicted \(K_d\), within the limits of the training set and the empirical model.
\end{itemize}

For virtual screening applications, absolute values of \(S_{\text{Surflex}}\) are typically less important than \emph{relative} scores across a congeneric series or a screening library.\cite{CH7_Jain2007Surflex21,CH7_Jain2006NegTrain}



\section{GOLD}

GOLD (Genetic Optimisation for Ligand Docking) is a widely used protein--ligand docking program that uses a genetic algorithm (GA) to explore ligand binding modes in a (partially) flexible protein binding site.\cite{CH7_Jones1997} It is developed and distributed by the Cambridge Crystallographic Data Centre (CCDC).

GOLD is particularly known for:

\begin{itemize}
    \item robust pose prediction with full ligand flexibility and optional side–chain and water flexibility;
    \item multiple scoring functions (GoldScore, ChemScore, ASP, ChemPLP) and consensus scoring;\cite{CH7_Verdonk2003,CH7_Mooij2005}
    \item explicit handling of binding-site water molecules, including displacement and reorientation;
    \item highly configurable constraints (distance, H-bond, pharmacophore, scaffold, etc.) to encode prior knowledge.
\end{itemize}

GOLD's design is typical of modern docking programs: a stochastic search procedure (here, a GA) generates candidate poses, and an empirical or knowledge-based scoring function evaluates their fitness.\cite{CH7_Jones1997} The combination of a flexible GA and several complementary scoring functions makes GOLD a good teaching platform for beginners.

\subsection{Fundamental working theory of GOLD}


In structure-based docking, the goal is to predict the bound pose of a ligand in a protein binding site, and, ideally, to estimate its binding affinity. Conceptually, one must:
\begin{enumerate}
    \item sample the relevant conformational and orientational space of the ligand (and sometimes the receptor);
    \item evaluate the quality of each sampled pose with a scoring function that approximates the binding free energy.
\end{enumerate}

The search space is high-dimensional: it includes ligand translations and rotations, rotatable bonds, ring conformations, and, in GOLD, possible flips of certain protein side chains and on/off states and orientations of some water molecules.\cite{CH7_Jones1997}

\subsubsection{Genetic algorithm in GOLD}

GOLD uses a steady-state, operator-based genetic algorithm to perform the conformational search.\cite{CH7_Jones1997} A genetic algorithm is an evolutionary optimisation method that maintains a population of candidate solutions (``chromosomes'') and iteratively improves them through selection, crossover, mutation, and (optionally) migration between subpopulations (islands).
\begin{enumerate}
\item Encoding (chromosomes)
In GOLD, each chromosome encodes a complete ligand pose:\cite{CH7_Jones1997}
\begin{itemize}
    \item the three translational degrees of freedom of the ligand;
    \item three rotational degrees of freedom (e.g.\ Euler angles or quaternions);
    \item values for each rotatable bond in the ligand;
    \item optional discrete variables for ring conformations;
    \item optional side-chain torsion states for flexible protein residues near the binding site;
    \item optional on/off (bound/displaced) flags and orientations for selected water molecules.
\end{itemize}
Each chromosome thus represents a complete candidate docking solution.

\item Fitness evaluation
The fitness (objective) function is a docking score: GoldScore, ChemScore, ASP, or ChemPLP. Higher scores correspond to better predicted binding (GOLD reports a dimensionless fitness function rather than an energy in physical units).

\item GA cycle.
The GA proceeds as follows:\cite{CH7_Jones1997}
\begin{enumerate}
    \item \textbf{Initialisation:} generate a random or biased initial population of chromosomes obeying simple steric constraints.
    \item \textbf{Evaluation:} compute the fitness score for each chromosome.
    \item \textbf{Selection:} select parent chromosomes with probability biased towards higher fitness (e.g.\ tournament or roulette-wheel selection).
    \item \textbf{Variation:} apply genetic operators:
    \begin{itemize}
        \item \emph{crossover} — recombine two parents to produce offspring (e.g.\ swapping subsets of torsions or pose parameters);
        \item \emph{mutation} — randomly perturb a degree of freedom (e.g.\ rotate one torsion, translate or rotate the ligand slightly, flip a side chain).
    \end{itemize}
    \item \textbf{Replacement (steady state):} insert offspring into the population, replacing some of the less fit chromosomes, often with duplicate suppression so that identical solutions are not preserved.\cite{CH7_Jones1997}
    \item \textbf{Termination:} after a fixed number of genetic operations or when convergence criteria (e.g.\ several best poses within a small RMSD) are met, the GA stops and returns the highest scoring poses.
\end{enumerate}
\end{enumerate}
GOLD uses an island model, where the population is partitioned into islands that evolve semi-independently with occasional migration. This encourages diversity and helps avoid premature convergence.\cite{CH7_Jones1997}

\subsection{Protein and water flexibility}

While the receptor is usually treated as rigid, GOLD supports limited protein and water flexibility:
\begin{itemize}
    \item \textbf{Protein side chains:} selected residues in the binding site can be treated as flexible, typically restricted to a library of side-chain rotamers or sampled torsions.
    \item \textbf{Water molecules:} explicitly modelled waters can be allowed to:
    \begin{itemize}
        \item toggle between bound (in the site) and displaced (returned to bulk);
        \item rotate to optimise hydrogen bonding;
        \item (optionally) translate within a small radius.
    \end{itemize}
\end{itemize}
GOLD estimates the free-energy change associated with bringing a water molecule from bulk into a particular site, and this contributes to the fitness of poses where that water is retained versus displaced.

\subsection{Typical GOLD workflow for beginners}

A standard GOLD workflow can be broken into the following stages.
\begin{enumerate}
  \item Protein (receptor) preparation

\begin{enumerate}
    \item \textbf{Start from a high-quality structure.} Use a crystal structure or cryo-EM structure with a bound ligand when possible.
    \item \textbf{Clean the structure:} remove alternate conformations, resolve missing atoms if possible, and inspect for crystallographic artefacts.
    \item \textbf{Add hydrogen atoms:} add all hydrogens and choose appropriate protonation/tautomeric states for titratable residues given the target pH.
    \item \textbf{Handle cofactors and metals:} retain essential metal ions and prosthetic groups; ensure their coordination geometry is reasonable.
    \item \textbf{Water molecules:} 
    \begin{itemize}
        \item keep structural waters that are known or suspected to be functionally important;
        \item remove obvious bulk waters;
        \item flag selected waters as potentially ``toggleable'' in GOLD for binding/displacement analysis.
    \end{itemize}
\end{enumerate}

\item Binding-site definition

The binding site can be defined in several ways:
\begin{itemize}
    \item \textbf{From a reference ligand:} specify a co-crystallised ligand; GOLD then defines a spherical binding site of user-defined radius (e.g.\ 8--12~\AA) around the ligand.
    \item \textbf{By residues:} specify a set of binding-site residues; GOLD uses their coordinates to define the pocket region.
    \item \textbf{By coordinates:} specify a point and radius (useful when docking into apo structures or homology models).
\end{itemize}
For training purposes, defining the site around a known ligand is recommended, as it provides a clear structural context and an RMSD reference for pose validation.

\item Ligand preparation

\begin{enumerate}
    \item \textbf{2D to 3D:} build ligands or import from databases and generate 3D structures.
    \item \textbf{Protonation and tautomers:} enumerate plausible protomers and tautomers at the working pH; many failures in docking arise from incorrect protonation states.
    \item \textbf{Conformation:} optionally pre-minimise ligands with a small-molecule force field; GOLD will sample conformations during docking but a reasonable starting structure can help.
    \item \textbf{Define rotatable bonds:} GOLD automatically detects rotatable bonds, but manual curation (e.g.\ ring amides, conjugated bonds) is sometimes necessary.
\end{enumerate}

\item Docking setup in GOLD / Hermes

\begin{itemize}
    \item \textbf{Choose the scoring function:} ChemPLP is the default and generally recommended for pose prediction and virtual screening.
    \item \textbf{GA parameters:} select:
    \begin{itemize}
        \item population size,
        \item number of genetic operations,
        \item selection pressure,
        \item operator weights (crossover, mutation, migration).
    \end{itemize}
    The search efficiency parameter in recent versions provides a convenient, ligand-size-aware way to scale the number of GA operations.
    \item \textbf{Number of GA runs per ligand:} for pose prediction, 10--30 independent runs per ligand are typical; for virtual screening, fewer runs may be used for speed.
    \item \textbf{Constraints (optional but powerful):}
    \begin{itemize}
        \item enforce key hydrogen bonds or metal interactions;
        \item require a pharmacophore pattern or a scaffold match;
        \item apply distance or shape similarity constraints relative to a reference ligand.
    \end{itemize}
\end{itemize}

\item Running docking and analysing results

After docking:
\begin{enumerate}
    \item \textbf{Rank poses} by the chosen fitness function.
    \item \textbf{Inspect top poses} visually for:
    \begin{itemize}
        \item agreement with known SAR and chemistry;
        \item reasonable geometry (no severe clashes, realistic torsions);
        \item preservation or rational modification of key interactions.
    \end{itemize}
    \item \textbf{RMSD analysis:} for systems with known co-crystal structures, compute RMSD between predicted and experimental poses (RMSD $\lesssim 2.0$~\AA\ is commonly considered a successful pose).\cite{CH7_Jones1997}
    \item \textbf{Rescoring / consensus scoring:} re-score docked poses with alternative scoring functions (e.g.\ dock with ChemPLP, rescore with GoldScore and ChemScore) to reduce scoring-function-specific bias.\cite{CH7_Verdonk2003}
\end{enumerate}
\end{enumerate}

\subsection{Scoring functions in GOLD}

GOLD provides four main scoring functions:
\begin{itemize}
    \item GoldScore (force-field-like, early default);
    \item ChemScore (empirical, trained against binding affinities);
    \item ASP (Astex Statistical Potential, knowledge-based);\cite{CH7_Mooij2005}
    \item ChemPLP (empirical, piecewise linear potential, current default).
\end{itemize}
All are dimensionless ``fitness'' scores; higher values indicate better predicted binding.

\subsubsection{GoldScore}

GoldScore is the original GOLD scoring function and is designed primarily for pose prediction.\cite{CH7_Jones1997} Conceptually, it is a force-field-inspired function consisting of:
\begin{itemize}
    \item an external van der Waals term between protein and ligand;
    \item an external hydrogen-bonding term between protein and ligand;
    \item an internal ligand strain term (van der Waals and torsional components).
\end{itemize}

In the implementation used in docking benchmarks, GoldScore is typically written as:\cite{CH7_Wang2003_Suppl}
\begin{equation}
S_{\mathrm{GoldScore}}
= 1.375\, S_{\mathrm{vdw,ext}}
+ S_{\mathrm{hb,ext}}
+ S_{\mathrm{int}},
\label{eq:GoldScore}
\end{equation}
where:
\begin{itemize}
    \item $S_{\mathrm{vdw,ext}}$ is the van der Waals steric interaction between protein and ligand (4--8 potential);\cite{CH7_Wang2003_Suppl}
    \item $S_{\mathrm{hb,ext}}$ is the hydrogen-bonding interaction between protein and ligand;
    \item $S_{\mathrm{int}}$ is an internal ligand strain term based on van der Waals and torsional contributions.
\end{itemize}
The empirical factor 1.375 partially compensates for the absence of an explicit hydrophobic term.\cite{CH7_Wang2003_Suppl}

The external van der Waals term is computed as a sum over all protein–ligand atom pairs $(i,j)$:
\begin{equation}
S_{\mathrm{vdw,ext}}
= \sum_{i \in L}\sum_{j \in P}
E^{\mathrm{vdw}}_{ij}(r_{ij}),
\end{equation}
where $r_{ij}$ is the interatomic distance. A 4--8 Lennard–Jones-like potential is used for the external van der Waals interactions:
\begin{equation}
E^{\mathrm{vdw}}_{ij}(r_{ij})
= \epsilon_{ij}
\left[
\left(\frac{r_{ij,0}}{r_{ij}}\right)^{8}
-
2\left(\frac{r_{ij,0}}{r_{ij}}\right)^{4}
\right],
\end{equation}
where:
\begin{itemize}
    \item $r_{ij,0}$ is a distance parameter related to the sum of van der Waals radii of atoms $i$ and $j$;
    \item $\epsilon_{ij}$ is an interaction well depth parameter depending on the atom types and interaction class (e.g.\ hydrogen-bonding vs.\ non-polar contact).
\end{itemize}

The hydrogen-bond term is a distance- and angle-dependent potential:\cite{CH7_Jones1997}
\begin{equation}
S_{\mathrm{hb,ext}}
=
\sum_{\mathrm{HB\ pairs}}
E^{\mathrm{hb}}_{ij}(r_{ij}, \theta_{ij}),
\end{equation}
where $E^{\mathrm{hb}}_{ij}$ is maximally favourable near an ideal donor–acceptor distance and angle, and decays as the geometry deviates from ideal hydrogen bonding.

Finally, the internal term $S_{\mathrm{int}}$ penalises high-energy ligand conformations (steric clashes, unfavourable torsions), discouraging poses that require unrealistic ligand strain.

\subsubsection{ChemScore}

ChemScore is an empirical scoring function originally developed by Eldridge \emph{et al.}\ for estimating protein–ligand binding free energies from structures.\cite{CH7_Eldridge1997} It was later implemented and adapted in GOLD.\cite{CH7_Verdonk2003} In its original form, ChemScore estimates the binding free energy as:
\begin{equation}
\Delta G_{\mathrm{bind}}
=
\Delta G_{0}
+
\Delta G_{\mathrm{hbond}}
\sum_{i} H_{i}
+
\Delta G_{\mathrm{metal}}
\sum_{j} M_{j}
+
\Delta G_{\mathrm{lipo}}
\sum_{k} L_{k}
+
\Delta G_{\mathrm{rot}} N_{\mathrm{rot}},
\label{eq:ChemScore}
\end{equation}
where:
\begin{itemize}
    \item $\Delta G_{0}$ is a constant offset (overall regression intercept);
    \item $\sum_{i} H_{i}$ is a sum over hydrogen-bonding interactions, each weighted by a geometry-dependent factor and the strength of the donor/acceptor;
    \item $\sum_{j} M_{j}$ is a sum over metal–ligand interactions (for ligands coordinating metal ions);
    \item $\sum_{k} L_{k}$ is a measure of lipophilic contact area (buried hydrophobic surface or hydrophobic contact count);
    \item $N_{\mathrm{rot}}$ is the number of rotatable bonds in the ligand in the unbound state;
    \item $\Delta G_{\mathrm{hbond}}$, $\Delta G_{\mathrm{metal}}$, $\Delta G_{\mathrm{lipo}}$, and $\Delta G_{\mathrm{rot}}$ are empirical coefficients obtained by regression against a training set of 82 complexes with known binding affinities.\cite{CH7_Eldridge1997}
\end{itemize}

In GOLD, the ChemScore-based fitness\cite{CH7_Wang2003_Suppl} incorporates additional penalties:
\begin{equation}
S_{\mathrm{ChemScore}}^{\mathrm{GOLD}}
=
-\Delta G_{\mathrm{bind}}
+
P_{\mathrm{clash}}
+
P_{\mathrm{torsion,int}},
\end{equation}
where:
\begin{itemize}
    \item $P_{\mathrm{clash}}$ penalises protein–ligand atom clashes;
    \item $P_{\mathrm{torsion,int}}$ is an internal torsion penalty discouraging energetically unfavourable ligand conformations (complementary to $N_{\mathrm{rot}}$).
\end{itemize}
A higher $S_{\mathrm{ChemScore}}^{\mathrm{GOLD}}$ corresponds to a more favourable (more negative) $\Delta G_{\mathrm{bind}}$ and fewer penalties.

\subsubsection{ASP: Astex Statistical Potential}

ASP is a knowledge-based, atom–atom distance potential derived from a database of protein–ligand complexes.\cite{CH7_Mooij2005} It is conceptually related to PMF and DrugScore potentials but uses an improved reference state that accounts for differences in exposure of protein atom types towards ligand binding sites.\cite{CH7_Mooij2005}

In general, a statistical potential for atom types $i$ and $j$ at separation $r$ can be written as:
\begin{equation}
w_{ij}(r)
=
- k_{\mathrm{B}} T
\ln
\left(
\frac{\rho_{ij}^{\mathrm{obs}}(r)}{\rho_{ij}^{\mathrm{ref}}(r)}
\right),
\label{eq:ASP}
\end{equation}
where:
\begin{itemize}
    \item $k_{\mathrm{B}}$ is Boltzmann's constant;
    \item $T$ is an effective temperature (often absorbed into the scale of the potential);
    \item $\rho_{ij}^{\mathrm{obs}}(r)$ is the observed distribution of $i$--$j$ distances in a structural database of complexes;
    \item $\rho_{ij}^{\mathrm{ref}}(r)$ is the reference (non-interacting) distribution that corrects for background atom exposure and other biases.
\end{itemize}

The ASP score for a pose is then:
\begin{equation}
S_{\mathrm{ASP}}
=
\sum_{i \in P}\sum_{j \in L}
w_{ij}(r_{ij})
+
S_{\mathrm{ChemScore,aux}},
\end{equation}
where $S_{\mathrm{ChemScore,aux}}$ denotes additional ChemScore-like terms (e.g.\ for hydrogen bonding and rotational entropy) that are incorporated into the ASP fitness in GOLD.

Because ASP is knowledge-based, it captures typical contact preferences (e.g.\ aromatic stacking, specific polar contacts) encoded in the database. It often performs comparably to GoldScore and ChemScore for pose prediction and binding-affinity ranking.\cite{CH7_Mooij2005}

\subsubsection{ChemPLP: piecewise linear potential}

ChemPLP is the current default scoring function in GOLD and is optimised for pose prediction and virtual screening.\ It combines:
\begin{itemize}
    \item a piecewise linear potential (PLP) for steric (van der Waals-like) protein–ligand interactions;
    \item ChemScore-like hydrogen-bond and metal-binding terms;
    \item internal ligand clash and torsion terms.
\end{itemize}

In schematic form, the ChemPLP score can be written as:
\begin{equation}
S_{\mathrm{ChemPLP}}
=
\sum_{i \in L}\sum_{j \in P}
f_{\mathrm{PLP}}(r_{ij})
+
\sum_{\mathrm{HB}}
f_{\mathrm{HB}}(r,\theta)
+
S_{\mathrm{lig,clash}}
+
S_{\mathrm{torsion}},
\label{eq:ChemPLP}
\end{equation}
where:
\begin{itemize}
    \item $f_{\mathrm{PLP}}(r_{ij})$ is a piecewise linear approximation to a steric potential, with different linear segments modelling strong repulsion at short range, optimal contact, and decay at longer range;
    \item $f_{\mathrm{HB}}(r,\theta)$ is a distance- and angle-dependent hydrogen-bond term derived from ChemScore;
    \item $S_{\mathrm{lig,clash}}$ penalises internal heavy-atom clashes within the ligand;
    \item $S_{\mathrm{torsion}}$ is a torsional strain term similar to that used in ChemScore.
\end{itemize}

The details of the breakpoints and slopes in $f_{\mathrm{PLP}}(r)$ are proprietary, but conceptually the PLP term rewards close, well-packed contacts while sharply penalising severe overlaps. Validation studies by CCDC and others have shown ChemPLP to provide the highest success rates for pose prediction among the GOLD scoring functions, which is why it is the default choice in GOLD~5.1 and later.

\subsection{Choosing and using scoring functions in training}

For teaching beginners, the following practical guidelines are useful:
\begin{itemize}
    \item Use \textbf{ChemPLP} as the default scoring function for pose prediction and small to medium virtual screening tasks.
    \item Demonstrate \textbf{GoldScore} to illustrate a more force-field-like potential and discuss its energy components explicitly.
    \item Use \textbf{ChemScore} when emphasising empirical free-energy decomposition into hydrogen bonding, lipophilic, metal, and rotational terms, and stress that absolute values are not true physical $\Delta G$.
    \item Present \textbf{ASP} to introduce knowledge-based potentials and how they are derived from structural statistics.
    \item Encourage \textbf{consensus scoring} (e.g.\ dock with ChemPLP, rescore top poses with GoldScore and ChemScore) to highlight the strengths and weaknesses of different scoring functions and to reduce dependence on any single model.\cite{CH7_Verdonk2003}
\end{itemize}

\section{Summary}

\bibliographystyle{plain}
\bibliography{Molecular_Docking_Methods_and_Programs}

\end{document}